% =====================================================================
\chapter{Soluciones de los ejercicios}
\label{ap:soluciones}
% =====================================================================

Cada solución se presenta en cinco partes: datos, planteamiento, desarrollo, resultado y comprobación. Los valores aproximados se dan con cuatro cifras significativas salvo que se indique otra cosa; los resultados exactos se dan en forma cerrada cuando existe. En los ejercicios de demostración, el apartado «Datos» recoge las hipótesis y el «Resultado», la conclusión.

% ---------------------------------------------------------------------
\section{Tema 1}

\begin{solucion}{eje:1-conversion}
\datos ángulos $135^{\circ}$, $210^{\circ}$, $-36^{\circ}$, $1^{\circ}$; $5\pi/6$, $7\pi/4$ y $2$ rad.
\planteamiento por la \cref{def:radian}, se multiplica por $\pi/180$ para pasar a radianes y por $180/\pi$ para pasar a grados.
\pasos
\begin{enumerate}[label=\arabic*.]
  \item $135\cdot\frac{\pi}{180}=\frac{3\pi}{4}$; \quad $210\cdot\frac{\pi}{180}=\frac{7\pi}{6}$; \quad $-36\cdot\frac{\pi}{180}=-\frac{\pi}{5}$; \quad $1\cdot\frac{\pi}{180}\approx0{,}017\,45$.
  \item $\frac{5\pi}{6}\cdot\frac{180}{\pi}=150^{\circ}$; \quad $\frac{7\pi}{4}\cdot\frac{180}{\pi}=315^{\circ}$; \quad $2\cdot\frac{180}{\pi}=\frac{360}{\pi}\approx114{,}6^{\circ}$.
\end{enumerate}
\resultado $\mathbf{3\pi/4}$, $\mathbf{7\pi/6}$, $\mathbf{-\pi/5}$, $\mathbf{\pi/180\approx0{,}01745}$ rad; $\mathbf{150^{\circ}}$, $\mathbf{315^{\circ}}$, $\mathbf{114{,}6^{\circ}}$.
\comprobacion $\frac{3\pi}{4}$ es tres cuartos de media vuelta, y $\frac34\cdot180^{\circ}=135^{\circ}$.
\end{solucion}

\begin{solucion}{eje:1-cuadrantes}
\datos $\theta\in\{2,4,5\}$; $\pi/2\approx1{,}57$, $\pi\approx3{,}14$, $3\pi/2\approx4{,}71$, $2\pi\approx6{,}28$.
\planteamiento situar cada $\theta$ entre los múltiplos de $\pi/2$ y aplicar el \cref{cor:signos}.
\pasos
\begin{enumerate}[label=\arabic*.]
  \item $1{,}57<2<3{,}14$: cuadrante II; $\sin2>0$, $\cos2<0$, $\tan2<0$.
  \item $3{,}14<4<4{,}71$: cuadrante III; $\sin4<0$, $\cos4<0$, $\tan4>0$.
  \item $4{,}71<5<6{,}28$: cuadrante IV; $\sin5<0$, $\cos5>0$, $\tan5<0$.
\end{enumerate}
\resultado \textbf{II $(+,-,-)$; III $(-,-,+)$; IV $(-,+,-)$} para $(\sin,\cos,\tan)$.
\comprobacion con calculadora en radianes: $\sin2=0{,}9093$, $\cos2=-0{,}4161$; $\sin4=-0{,}7568$, $\cos4=-0{,}6536$; $\sin5=-0{,}9589$, $\cos5=0{,}2837$.
\end{solucion}

\begin{solucion}{eje:1-vf}
\planteamiento buscar un contraejemplo o una justificación general para cada afirmación.
\pasos
\begin{enumerate}[label=(\alph*)]
  \item \textbf{Falsa.} Con $x=\pi/2$: $\sin^{2}(\pi/2)=1$, mientras que $\sin(\pi^{2}/4)\approx\sin2{,}467\approx0{,}6243$.
  \item \textbf{Falsa.} $\sin^{-1}$ denota la función inversa: $\sin^{-1}\frac12=\arcsin\frac12=\frac{\pi}{6}\approx0{,}5236$, mientras que $1/\sin\frac12\approx2{,}086$.
  \item \textbf{Falsa.} Por el \cref{teo:simetrias}(i), $\abs{\cos\theta}\le1$.
  \item \textbf{Verdadera.} Si $\theta-\varphi\in2\pi\Z$, el \cref{teo:simetrias}(ii) da $P(\theta)=P(\varphi)$: mismo seno y mismo coseno, luego misma tangente.
\end{enumerate}
\resultado \textbf{(a) falsa; (b) falsa; (c) falsa; (d) verdadera.}
\end{solucion}

\begin{solucion}{eje:1-arco}
\datos $r=6$ cm; $\theta=150^{\circ}$; arco de $9$ cm.
\planteamiento \cref{prop:arco-sector}: $s=r\theta$ y $A=\frac12r^{2}\theta$, con $\theta$ en radianes.
\pasos
\begin{enumerate}[label=\arabic*.]
  \item $\theta=150\cdot\frac{\pi}{180}=\frac{5\pi}{6}$.
  \item $s=6\cdot\frac{5\pi}{6}=5\pi\approx15{,}71$ cm.
  \item $A=\frac12\cdot36\cdot\frac{5\pi}{6}=15\pi\approx47{,}12$ cm$^{2}$.
  \item Para $s=9$: $\theta=s/r=9/6=1{,}5$ rad $=1{,}5\cdot\frac{180}{\pi}\approx85{,}94^{\circ}$.
\end{enumerate}
\resultado \textbf{$s=5\pi\approx15{,}71$ cm, $A=15\pi\approx47{,}12$ cm$^{2}$; el arco de 9 cm corresponde a $1{,}5$ rad $\approx85{,}94^{\circ}$.}
\comprobacion $150^{\circ}$ es $\frac{5}{12}$ de vuelta: $\frac{5}{12}\cdot2\pi\cdot6=5\pi$ y $\frac{5}{12}\cdot36\pi=15\pi$.
\end{solucion}

\begin{solucion}{eje:1-exactos}
\planteamiento para cada ángulo, cuadrante, ángulo de referencia (\cref{tab:notables}) y signo (\cref{cor:signos}).
\pasos
\begin{enumerate}[label=\arabic*.]
  \item $\frac{5\pi}{6}$: cuadrante II, referencia $\frac{\pi}{6}$; $\sin\frac{5\pi}{6}=+\frac12$.
  \item $\frac{7\pi}{4}$: cuadrante IV, referencia $\frac{\pi}{4}$; $\cos\frac{7\pi}{4}=+\frac{\sqrt2}{2}$.
  \item $\frac{4\pi}{3}$: cuadrante III, referencia $\frac{\pi}{3}$; tangente positiva: $\tan\frac{4\pi}{3}=\sqrt3$.
  \item $-\frac{\pi}{3}$: cuadrante IV, referencia $\frac{\pi}{3}$; $\sin(-\frac{\pi}{3})=-\frac{\sqrt3}{2}$.
  \item $\frac{11\pi}{6}$: cuadrante IV, referencia $\frac{\pi}{6}$; $\cos\frac{11\pi}{6}=+\frac{\sqrt3}{2}$.
  \item $\frac{3\pi}{2}$: sobre el eje, $P(\frac{3\pi}{2})=(0,-1)$; $\sin\frac{3\pi}{2}=-1$.
\end{enumerate}
\resultado $\mathbf{\tfrac12,\ \tfrac{\sqrt2}{2},\ \sqrt3,\ -\tfrac{\sqrt3}{2},\ \tfrac{\sqrt3}{2},\ -1}$.
\comprobacion $\frac{4\pi}{3}=\pi+\frac{\pi}{3}$, y el \cref{teo:simetrias}(v) da $\tan(\frac{\pi}{3}+\pi)=\frac{-\sin(\pi/3)}{-\cos(\pi/3)}=\tan\frac{\pi}{3}=\sqrt3$.
\end{solucion}

\begin{solucion}{eje:1-seis-razones}
\datos $\sin\theta=3/5$; $\theta$ en el cuadrante II.
\planteamiento identidad pitagórica y signo del coseno en el cuadrante II.
\pasos
\begin{enumerate}[label=\arabic*.]
  \item $\cos^{2}\theta=1-\frac{9}{25}=\frac{16}{25}$; en el cuadrante II, $\cos\theta=-\frac45$.
  \item $\tan\theta=\frac{3/5}{-4/5}=-\frac34$.
  \item $\sec\theta=-\frac54$, $\csc\theta=\frac53$, $\cot\theta=-\frac43$.
\end{enumerate}
\resultado $\mathbf{\cos\theta=-\tfrac45,\ \tan\theta=-\tfrac34,\ \sec\theta=-\tfrac54,\ \csc\theta=\tfrac53,\ \cot\theta=-\tfrac43}$.
\comprobacion $1+\tan^{2}\theta=1+\frac{9}{16}=\frac{25}{16}=\sec^{2}\theta$ (\cref{prop:pitagoricas}).
\end{solucion}

\begin{solucion}{eje:1-elevacion}
\datos cateto contiguo $30$ m; ángulo $60^{\circ}$.
\planteamiento \cref{def:razones-triangulo}: altura $h=30\tan60^{\circ}$; distancia $d=30/\cos60^{\circ}$.
\pasos
\begin{enumerate}[label=\arabic*.]
  \item $h=30\sqrt3\approx51{,}96$ m.
  \item $d=30/\frac12=60$ m.
\end{enumerate}
\resultado \textbf{altura $30\sqrt3\approx51{,}96$ m; distancia a la parte alta $60$ m.}
\comprobacion Pitágoras: $30^{2}+(30\sqrt3)^{2}=900+2700=3600=60^{2}$.
\end{solucion}

\begin{solucion}{eje:1-coterminales}
\datos $1020^{\circ}$, $-11\pi/4$, $17\pi/6$.
\planteamiento sumar o restar vueltas completas hasta caer en $[0^{\circ},360^{\circ})$ o $[0,2\pi)$; después, cuadrante y referencia.
\pasos
\begin{enumerate}[label=\arabic*.]
  \item $1020^{\circ}-2\cdot360^{\circ}=300^{\circ}$: cuadrante IV, referencia $60^{\circ}$; $\sin=-\frac{\sqrt3}{2}$, $\cos=\frac12$.
  \item $-\frac{11\pi}{4}+4\pi=\frac{5\pi}{4}$: cuadrante III, referencia $\frac{\pi}{4}$; $\sin=\cos=-\frac{\sqrt2}{2}$.
  \item $\frac{17\pi}{6}-2\pi=\frac{5\pi}{6}$: cuadrante II, referencia $\frac{\pi}{6}$; $\sin=\frac12$, $\cos=-\frac{\sqrt3}{2}$.
\end{enumerate}
\resultado \textbf{$300^{\circ}$ $(-\frac{\sqrt3}{2},\frac12)$; $\frac{5\pi}{4}$ $(-\frac{\sqrt2}{2},-\frac{\sqrt2}{2})$; $\frac{5\pi}{6}$ $(\frac12,-\frac{\sqrt3}{2})$}, dados como $(\sin,\cos)$.
\comprobacion en cada caso $\sin^{2}+\cos^{2}=1$: por ejemplo, $\frac34+\frac14=1$.
\end{solucion}

\begin{solucion}{eje:1-rope}
\datos ángulo $m$ rad en la posición $m=100$.
\planteamiento dividir entre $2\pi$; el resto es el ángulo coterminal. Para (c), usar el \cref{teo:simetrias}(ii).
\pasos
\begin{enumerate}[label=\arabic*.]
  \item $100/(2\pi)\approx15{,}92$: se han dado $15$ vueltas completas, y $100-15\cdot2\pi=100-30\pi\approx100-94{,}25=5{,}752$ rad ($\approx329{,}6^{\circ}$).
  \item $3\pi/2\approx4{,}712<5{,}752<2\pi\approx6{,}283$: cuadrante IV, seno negativo y coseno positivo.
  \item Las posiciones $m\ne m'$ darían exactamente el mismo punto si y solo si $m-m'=2\pi k$ con $k\in\Z$ no nulo, es decir, si $\pi=(m-m')/(2k)$ fuera racional. Como $\pi$ es irracional, eso no ocurre nunca.
\end{enumerate}
\resultado \textbf{(a) 15 vueltas, ángulo $100-30\pi\approx5{,}752$ rad; (b) cuadrante IV, $\sin<0$, $\cos>0$; (c) nunca se repite exactamente}, aunque se aproxima a una vuelta cada $2\pi\approx6{,}28$ posiciones.
\comprobacion $\sin100\approx-0{,}5064$ y $\cos100\approx0{,}8623$.
\end{solucion}

\begin{solucion}{eje:1-simplificar}
\planteamiento \cref{teo:simetrias}, sin desarrollar fórmulas de la suma.
\pasos
\begin{enumerate}[label=(\alph*)]
  \item $\sin(\pi-x)=\sin x$; $\cos(\frac{\pi}{2}+x)=\cos(\frac{\pi}{2}-(-x))=\sin(-x)=-\sin x$. Suma: $0$.
  \item $\cos(x+\pi)\cos(-x)=-\cos x\cdot\cos x=-\cos^{2}x$; $\sin(x+\pi)\sin(\pi-x)=-\sin x\cdot\sin x=-\sin^{2}x$. Diferencia: $-\cos^{2}x+\sin^{2}x$.
  \item $\tan(\pi+x)=\frac{-\sin x}{-\cos x}=\tan x$; $\cot(\pi-x)=\frac{-\cos x}{\sin x}=-\cot x$. Producto: $-\tan x\cot x=-1$.
\end{enumerate}
\resultado \textbf{(a) $0$; (b) $\sin^{2}x-\cos^{2}x$; (c) $-1$.}
\comprobacion (b) con $x=0$: $\cos\pi\cos0-\sin\pi\sin\pi=-1=\sin^{2}0-\cos^{2}0$.
\end{solucion}

\begin{solucion}{eje:1-sinx-x-tanx}
\datos $0<x<\pi/2$.
\planteamiento con $O=(0,0)$, $A=(1,0)$, $P=P(x)=(\cos x,\sin x)$ y $T=(1,\tan x)$, punto en que la recta $OP$ corta a la recta $x=1$, se comparan tres áreas encajadas.
\pasos
\begin{enumerate}[label=\arabic*.]
  \item El triángulo $OAP$ tiene base $OA=1$ y altura $\sin x$: área $\frac12\sin x$.
  \item El sector circular $OAP$ tiene radio~1 y ángulo $x$: área $\frac12x$ (\cref{prop:arco-sector}).
  \item El triángulo $OAT$ es rectángulo en $A$, con catetos $1$ y $\tan x$: área $\frac12\tan x$.
  \item El triángulo $OAP$ está contenido en el sector (la cuerda $AP$ queda dentro del círculo), y el sector en el triángulo $OAT$ (el arco queda por debajo de la recta tangente $x=1$ y por debajo de la recta $OT$). Las inclusiones son estrictas, porque las diferencias contienen regiones de área positiva.
\end{enumerate}
\resultado $\frac12\sin x<\frac12x<\frac12\tan x$, es decir, $\mathbf{\sin x<x<\tan x}$.
\comprobacion para $x=\pi/6$: $0{,}5<0{,}5236<0{,}5774$.
\end{solucion}

\begin{solucion}{eje:1-abs-sin}
\datos \cref{eje:1-sinx-x-tanx}; $\abs{\sin x}\le1$; paridad del seno.
\planteamiento distinguir $x=0$, $0<x<\pi/2$, $x\ge\pi/2$ y usar la paridad para $x<0$.
\pasos
\begin{enumerate}[label=\arabic*.]
  \item $x=0$: igualdad, $0=0$.
  \item $0<x<\pi/2$: $0<\sin x<x$ por el \cref{eje:1-sinx-x-tanx}, luego $\abs{\sin x}<\abs{x}$.
  \item $x\ge\pi/2$: $\abs{\sin x}\le1<\pi/2\le x$, porque $\pi/2\approx1{,}571>1$.
  \item $x<0$: $\abs{\sin x}=\abs{\sin(-x)}<\abs{-x}=\abs{x}$, por los dos pasos anteriores aplicados a $-x>0$.
\end{enumerate}
\resultado \textbf{$\abs{\sin x}\le\abs{x}$ para todo $x$, con igualdad solo en $x=0$.}
\comprobacion $\abs{\sin(-2)}=0{,}9093<2$.
\end{solucion}

\begin{solucion}{eje:1-sin-mas-cos}
\datos $s=\sin\theta$, $c=\cos\theta$, $s^{2}+c^{2}=1$.
\planteamiento acotar $(s+c)^{2}$ usando $(s-c)^{2}\ge0$.
\pasos
\begin{enumerate}[label=\arabic*.]
  \item $(s-c)^{2}\ge0$ da $2sc\le s^{2}+c^{2}=1$.
  \item $(s+c)^{2}=s^{2}+c^{2}+2sc=1+2sc\le2$, luego $\abs{s+c}\le\sqrt2$.
  \item Hay igualdad $s+c=\sqrt2$ si y solo si $(s-c)^{2}=0$ y $s+c>0$, es decir, $s=c=\frac{\sqrt2}{2}$, lo que ocurre si y solo si $P(\theta)=P(\pi/4)$.
\end{enumerate}
\resultado \textbf{$\abs{\sin\theta+\cos\theta}\le\sqrt2$, con $\sin\theta+\cos\theta=\sqrt2$ si y solo si $\theta=\frac{\pi}{4}+2k\pi$.}
\comprobacion la \cref{prop:armonico} da $\sin\theta+\cos\theta=\sqrt2\sin(\theta+\frac{\pi}{4})$, cuyo máximo es $\sqrt2$ en $\theta+\frac{\pi}{4}=\frac{\pi}{2}+2k\pi$.
\end{solucion}

\begin{solucion}{eje:1-periodo-tan}
\datos \cref{teo:simetrias}; $\tan=\sin/\cos$ donde $\cos\ne0$.
\planteamiento (a) los puntos de la circunferencia con ordenada~0; (b) una contradicción con (a).
\pasos
\begin{enumerate}[label=\arabic*.]
  \item (a) $\sin\theta=0$ si y solo si $P(\theta)\in\{(1,0),(-1,0)\}=\{P(0),P(\pi)\}$, es decir, por el \cref{teo:simetrias}(ii), si $\theta\in2\pi\Z$ o $\theta\in\pi+2\pi\Z$. La unión de ambos conjuntos es $\pi\Z$.
  \item (b) Si $\cos\theta\ne0$, entonces $\cos(\theta+\pi)=-\cos\theta\ne0$ y $\tan(\theta+\pi)=\frac{-\sin\theta}{-\cos\theta}=\tan\theta$.
  \item Supóngase que $0<T<\pi$ cumple $\tan(\theta+T)=\tan\theta$ en todo el dominio. Con $\theta=0$: $\tan T=\tan0=0$, luego $\sin T=0$ y, por (a), $T\in\pi\Z$, imposible si $0<T<\pi$.
\end{enumerate}
\resultado \textbf{$\sin\theta=0\iff\theta\in\pi\Z$, y el periodo fundamental de la tangente es $\pi$.}
\comprobacion $\tan\frac{\pi}{4}=1$ y $\tan\frac{5\pi}{4}=1$, mientras que $\tan(\frac{\pi}{4}+\frac{\pi}{2})=\tan\frac{3\pi}{4}=-1$: $\pi/2$ no es periodo.
\end{solucion}

\begin{solucion}{eje:1-poligono}
\datos $n\ge3$; vértices $P(2\pi k/n)$, $k=0,\dots,n-1$.
\planteamiento calcular un lado con el triángulo isósceles que forma con el centro, y acotar con el \cref{eje:1-sinx-x-tanx}.
\pasos
\begin{enumerate}[label=\arabic*.]
  \item Cada lado y el centro forman un triángulo isósceles con dos lados iguales a~1 y ángulo central $2\pi/n$. La altura desde el centro lo divide en dos triángulos rectángulos de hipotenusa~1 y ángulo $\pi/n$, cuyo cateto opuesto mide $\sin(\pi/n)$.
  \item Lado $=2\sin(\pi/n)$; perímetro $p_n=2n\sin(\pi/n)$.
  \item Como $n\ge3$, $0<\pi/n<\pi/2$, y $\sin(\pi/n)<\pi/n$; por tanto $p_n<2n\cdot\pi/n=2\pi$.
  \item $p_6=12\sin\frac{\pi}{6}=6$.
\end{enumerate}
\resultado \textbf{$p_n=2n\sin(\pi/n)<2\pi$; $p_6=6$.}
\comprobacion el hexágono regular inscrito tiene lado igual al radio, $1$, y perímetro $6<2\pi\approx6{,}283$. Además $p_{96}\approx6{,}282$, la aproximación de Arquímedes.
\end{solucion}

% ---------------------------------------------------------------------
\section{Tema 2}

\begin{solucion}{eje:2-vf}
\planteamiento contrastar con el \cref{cor:suma} y la \cref{prop:doble}; buscar contraejemplos.
\pasos
\begin{enumerate}[label=(\alph*)]
  \item \textbf{Falsa.} Con $\alpha=\beta=0$: $\cos0=1\ne\cos0+\cos0=2$.
  \item \textbf{Verdadera.} Es la \cref{prop:doble}.
  \item \textbf{Falsa.} Lo correcto es $2\cos^{2}\alpha-1$; con $\alpha=0$, $1-2\cos^{2}0=-1\ne\cos0=1$.
  \item \textbf{Verdadera.} \Cref{eje:1-periodo-tan}(b).
\end{enumerate}
\resultado \textbf{(a) falsa; (b) verdadera; (c) falsa; (d) verdadera.}
\end{solucion}

\begin{solucion}{eje:2-signo-mitad}
\datos $\pi<\alpha<2\pi$.
\planteamiento el signo lo fija el cuadrante de $\alpha/2$ (\cref{cor:mitad}).
\pasos
\begin{enumerate}[label=\arabic*.]
  \item Dividiendo por~2: $\pi/2<\alpha/2<\pi$, cuadrante II.
  \item En el cuadrante II el coseno es negativo y el seno positivo.
\end{enumerate}
\resultado $\mathbf{\cos\frac{\alpha}{2}=-\sqrt{(1+\cos\alpha)/2}}$ y $\mathbf{\sin\frac{\alpha}{2}=+\sqrt{(1-\cos\alpha)/2}}$.
\comprobacion con $\alpha=\frac{4\pi}{3}$: $\cos\frac{2\pi}{3}=-\frac12$ y $-\sqrt{(1-\frac12)/2}=-\frac12$.
\end{solucion}

\begin{solucion}{eje:2-exactos}
\planteamiento $75^{\circ}=45^{\circ}+30^{\circ}$ y $105^{\circ}=60^{\circ}+45^{\circ}$; \cref{cor:suma}.
\pasos
\begin{enumerate}[label=\arabic*.]
  \item $\sin75^{\circ}=\frac{\sqrt2}{2}\cdot\frac{\sqrt3}{2}+\frac{\sqrt2}{2}\cdot\frac12=\frac{\sqrt6+\sqrt2}{4}$.
  \item $\cos105^{\circ}=\frac12\cdot\frac{\sqrt2}{2}-\frac{\sqrt3}{2}\cdot\frac{\sqrt2}{2}=\frac{\sqrt2-\sqrt6}{4}$.
  \item $\tan105^{\circ}=\frac{\sqrt3+1}{1-\sqrt3}$; multiplicando arriba y abajo por $1+\sqrt3$: $\frac{(1+\sqrt3)^{2}}{1-3}=\frac{4+2\sqrt3}{-2}=-2-\sqrt3$.
\end{enumerate}
\resultado $\mathbf{\sin75^{\circ}=\frac{\sqrt6+\sqrt2}{4}}$, $\mathbf{\cos105^{\circ}=\frac{\sqrt2-\sqrt6}{4}}$ y $\mathbf{\tan105^{\circ}=-2-\sqrt3}$, es decir, $0{,}9659$, $-0{,}2588$ y $-3{,}732$.
\comprobacion $\sin75^{\circ}=\cos15^{\circ}$, que coincide con el \cref{ej:cos15}; y $\cos105^{\circ}=-\sin15^{\circ}$.
\end{solucion}

\begin{solucion}{eje:2-cuadrantes}
\datos $\sin\alpha=\frac35$ (I); $\cos\beta=-\frac{5}{13}$ (II).
\planteamiento completar con la identidad pitagórica y aplicar el \cref{cor:suma}.
\pasos
\begin{enumerate}[label=\arabic*.]
  \item $\cos\alpha=+\frac45$; $\sin\beta=+\sqrt{1-\frac{25}{169}}=\frac{12}{13}$.
  \item $\sin(\alpha+\beta)=\frac35\cdot(-\frac{5}{13})+\frac45\cdot\frac{12}{13}=\frac{-15+48}{65}=\frac{33}{65}$.
  \item $\cos(\alpha+\beta)=\frac45\cdot(-\frac{5}{13})-\frac35\cdot\frac{12}{13}=\frac{-20-36}{65}=-\frac{56}{65}$.
  \item $\tan(\alpha+\beta)=-\frac{33}{56}$; seno positivo y coseno negativo: cuadrante II.
\end{enumerate}
\resultado $\mathbf{\sin(\alpha+\beta)=\frac{33}{65}}$, $\mathbf{\cos(\alpha+\beta)=-\frac{56}{65}}$, $\mathbf{\tan(\alpha+\beta)=-\frac{33}{56}}$; \textbf{cuadrante II}.
\comprobacion $33^{2}+56^{2}=1089+3136=4225=65^{2}$.
\end{solucion}

\begin{solucion}{eje:2-doble}
\datos $\cos\theta=-\frac35$; $\theta$ en el cuadrante III.
\planteamiento obtener $\sin\theta$ y aplicar la \cref{prop:doble}.
\pasos
\begin{enumerate}[label=\arabic*.]
  \item En el cuadrante III, $\sin\theta=-\frac45$.
  \item $\sin2\theta=2\cdot(-\frac45)(-\frac35)=\frac{24}{25}$.
  \item $\cos2\theta=\frac{9}{25}-\frac{16}{25}=-\frac{7}{25}$.
  \item $\tan2\theta=\frac{24/25}{-7/25}=-\frac{24}{7}$.
\end{enumerate}
\resultado $\mathbf{\sin2\theta=\frac{24}{25},\ \cos2\theta=-\frac{7}{25},\ \tan2\theta=-\frac{24}{7}}$.
\comprobacion $24^{2}+7^{2}=625=25^{2}$; además, con $\tan\theta=\frac43$, $\frac{2\tan\theta}{1-\tan^{2}\theta}=\frac{8/3}{-7/9}=-\frac{24}{7}$.
\end{solucion}

\begin{solucion}{eje:2-pi8}
\datos $\alpha=\pi/4$, $\cos\alpha=\sin\alpha=\frac{\sqrt2}{2}$; $\pi/8$ en el cuadrante I.
\planteamiento \cref{cor:mitad} con signo positivo.
\pasos
\begin{enumerate}[label=\arabic*.]
  \item $\sin\frac{\pi}{8}=\sqrt{\frac{1-\sqrt2/2}{2}}=\sqrt{\frac{2-\sqrt2}{4}}=\frac{\sqrt{2-\sqrt2}}{2}$.
  \item $\cos\frac{\pi}{8}=\frac{\sqrt{2+\sqrt2}}{2}$, del mismo modo.
  \item $\tan\frac{\pi}{8}=\frac{1-\cos(\pi/4)}{\sin(\pi/4)}=\frac{1-\sqrt2/2}{\sqrt2/2}=\frac{2-\sqrt2}{\sqrt2}=\sqrt2-1$.
\end{enumerate}
\resultado $\mathbf{\sin\frac{\pi}{8}=\frac{\sqrt{2-\sqrt2}}{2}\approx0{,}3827}$, $\mathbf{\cos\frac{\pi}{8}=\frac{\sqrt{2+\sqrt2}}{2}\approx0{,}9239}$, $\mathbf{\tan\frac{\pi}{8}=\sqrt2-1\approx0{,}4142}$.
\comprobacion $0{,}3827/0{,}9239=0{,}4142$, y $\tan2\cdot\frac{\pi}{8}=\frac{2(\sqrt2-1)}{1-(3-2\sqrt2)}=\frac{2(\sqrt2-1)}{2(\sqrt2-1)}=1=\tan\frac{\pi}{4}$.
\end{solucion}

\begin{solucion}{eje:2-producto}
\planteamiento \cref{prop:producto-suma} y \cref{cor:suma-producto}.
\pasos
\begin{enumerate}[label=\arabic*.]
  \item $\cos3x\cos x=\frac12[\cos2x+\cos4x]$.
  \item $\sin4x\cos2x=\frac12[\sin6x+\sin2x]$.
  \item $\sin5x-\sin3x=2\cos\frac{8x}{2}\sin\frac{2x}{2}=2\cos4x\sin x$.
  \item $\cos x+\cos3x=2\cos\frac{4x}{2}\cos\frac{-2x}{2}=2\cos2x\cos x$.
\end{enumerate}
\resultado $\mathbf{\tfrac12(\cos2x+\cos4x)}$, $\mathbf{\tfrac12(\sin6x+\sin2x)}$, $\mathbf{2\cos4x\sin x}$, $\mathbf{2\cos2x\cos x}$.
\comprobacion con $x=\pi/6$ en el tercero: $\sin\frac{5\pi}{6}-\sin\frac{\pi}{2}=\frac12-1=-\frac12$ y $2\cos\frac{2\pi}{3}\sin\frac{\pi}{6}=2\cdot(-\frac12)\cdot\frac12=-\frac12$.
\end{solucion}

\begin{solucion}{eje:2-potencias}
\planteamiento aplicar dos veces la reducción de potencias (\cref{cor:mitad}).
\pasos
\begin{enumerate}[label=\arabic*.]
  \item $\sin^{4}x=\bigl(\frac{1-\cos2x}{2}\bigr)^{2}=\frac{1-2\cos2x+\cos^{2}2x}{4}$.
  \item $\cos^{2}2x=\frac{1+\cos4x}{2}$.
  \item $\sin^{4}x=\frac{1}{4}\bigl(\frac32-2\cos2x+\frac12\cos4x\bigr)=\frac38-\frac12\cos2x+\frac18\cos4x$.
\end{enumerate}
\resultado $\mathbf{\sin^{4}x=\frac38-\frac12\cos2x+\frac18\cos4x}$.
\comprobacion $x=0$: $\frac38-\frac12+\frac18=0$; $x=\frac{\pi}{2}$: $\frac38+\frac12+\frac18=1$. Con $x=0{,}7$ ambos miembros valen $0{,}1722$.
\end{solucion}

\begin{solucion}{eje:2-verificar}
\planteamiento transformar el miembro izquierdo; anotar dónde están definidos ambos miembros.
\pasos
\begin{enumerate}[label=\arabic*.]
  \item (a) $\frac{1-\cos2x}{\sin2x}=\frac{2\sin^{2}x}{2\sin x\cos x}=\frac{\sin x}{\cos x}=\tan x$. Ambos miembros están definidos si $\sin2x\ne0$, es decir, si $x\ne k\pi/2$.
  \item (b) Con denominador común $(1+\cos x)\sin x$, el numerador es $\sin^{2}x+(1+\cos x)^{2}=\sin^{2}x+1+2\cos x+\cos^{2}x=2(1+\cos x)$.
  \item Simplificando $1+\cos x$: $\frac{2}{\sin x}=2\csc x$. Hace falta $\sin x\ne0$ (lo que ya implica $\cos x\ne-1$), es decir, $x\ne k\pi$.
\end{enumerate}
\resultado \textbf{(a) válida para $x\ne k\pi/2$; (b) válida para $x\ne k\pi$.}
\comprobacion (b) con $x=\pi/2$: $\frac11+\frac11=2=2\csc\frac{\pi}{2}$.
\end{solucion}

\begin{solucion}{eje:2-triple}
\planteamiento $3x=2x+x$, \cref{cor:suma} y \cref{prop:doble}.
\pasos
\begin{enumerate}[label=\arabic*.]
  \item $\cos3x=\cos2x\cos x-\sin2x\sin x=(2\cos^{2}x-1)\cos x-2\sin^{2}x\cos x$.
  \item Con $\sin^{2}x=1-\cos^{2}x$: $2\cos^{3}x-\cos x-2\cos x+2\cos^{3}x=4\cos^{3}x-3\cos x$.
  \item $\sin3x=\sin2x\cos x+\cos2x\sin x=2\sin x\cos^{2}x+(1-2\sin^{2}x)\sin x$.
  \item Con $\cos^{2}x=1-\sin^{2}x$: $2\sin x-2\sin^{3}x+\sin x-2\sin^{3}x=3\sin x-4\sin^{3}x$.
\end{enumerate}
\resultado $\mathbf{\cos3x=4\cos^{3}x-3\cos x}$ y $\mathbf{\sin3x=3\sin x-4\sin^{3}x}$.
\comprobacion $x=\pi/3$: $\cos\pi=-1$ y $4\cdot\frac18-\frac32=-1$; $\sin\pi=0$ y $3\cdot\frac{\sqrt3}{2}-4\cdot\frac{3\sqrt3}{8}=0$.
\end{solucion}

\begin{solucion}{eje:2-cos36}
\datos $\theta=\pi/5$, de modo que $3\theta=\pi-2\theta$.
\planteamiento igualar $\sin3\theta=\sin2\theta$ y resolver la ecuación de segundo grado en $c=\cos\theta$.
\pasos
\begin{enumerate}[label=\arabic*.]
  \item Por el \cref{teo:simetrias}(iv), $\sin3\theta=\sin(\pi-2\theta)=\sin2\theta$.
  \item \Cref{eje:2-triple} y \cref{prop:doble}: $3\sin\theta-4\sin^{3}\theta=2\sin\theta\cos\theta$.
  \item Como $\sin\theta\ne0$, se divide: $3-4\sin^{2}\theta=2c$, y con $\sin^{2}\theta=1-c^{2}$: $4c^{2}-2c-1=0$.
  \item $c=\frac{2\pm\sqrt{4+16}}{8}=\frac{1\pm\sqrt5}{4}$. Como $\pi/5$ está en el cuadrante I, $c>0$, y $\frac{1-\sqrt5}{4}<0$ se descarta.
\end{enumerate}
\resultado $\mathbf{\cos\frac{\pi}{5}=\frac{1+\sqrt5}{4}\approx0{,}8090}$, la mitad de la razón áurea.
\comprobacion $\cos36^{\circ}=0{,}809\,02$ con calculadora, y $4(0{,}809)^{2}-2(0{,}809)-1\approx0$.
\end{solucion}

\begin{solucion}{eje:2-suma-senos}
\datos $\sin(x/2)\ne0$; $n\ge1$.
\planteamiento multiplicar cada término por $2\sin(x/2)$ y usar $2\sin\alpha\sin\beta=\cos(\alpha-\beta)-\cos(\alpha+\beta)$.
\pasos
\begin{enumerate}[label=\arabic*.]
  \item Con $\alpha=kx$ y $\beta=x/2$: $2\sin\frac{x}{2}\sin kx=\cos\bigl((k-\tfrac12)x\bigr)-\cos\bigl((k+\tfrac12)x\bigr)$.
  \item Sumando de $k=1$ a $n$, los términos intermedios se cancelan: $2\sin\frac{x}{2}\sum_{k=1}^{n}\sin kx=\cos\frac{x}{2}-\cos\bigl((n+\tfrac12)x\bigr)$.
  \item Por el \cref{cor:suma-producto}, $\cos A-\cos B=-2\sin\frac{A+B}{2}\sin\frac{A-B}{2}$ con $A=\frac x2$, $B=(n+\frac12)x$: $\frac{A+B}{2}=\frac{(n+1)x}{2}$ y $\frac{A-B}{2}=-\frac{nx}{2}$, luego la diferencia vale $2\sin\frac{(n+1)x}{2}\sin\frac{nx}{2}$.
  \item Se divide por $2\sin\frac x2\ne0$.
\end{enumerate}
\resultado $\displaystyle\sum_{k=1}^{n}\sin kx=\frac{\sin\frac{nx}{2}\sin\frac{(n+1)x}{2}}{\sin\frac x2}$.
\comprobacion $n=1$: $\frac{\sin\frac x2\sin x}{\sin\frac x2}=\sin x$. Con $n=7$, $x=0{,}9$, ambos miembros valen $0{,}008\,553$.
\end{solucion}

\begin{solucion}{eje:2-tan-triangulo}
\datos $\alpha+\beta+\gamma=\pi$; ningún ángulo recto.
\planteamiento $\gamma=\pi-(\alpha+\beta)$ y fórmula de la tangente de una suma.
\pasos
\begin{enumerate}[label=\arabic*.]
  \item Las tres tangentes están definidas porque ningún ángulo es $\pi/2$.
  \item $1-\tan\alpha\tan\beta\ne0$: si fuera cero, $\cos(\alpha+\beta)=\cos\alpha\cos\beta(1-\tan\alpha\tan\beta)=0$, y $\alpha+\beta=\pi/2$, con lo que $\gamma=\pi/2$, excluido.
  \item $\tan\gamma=\tan(\pi-(\alpha+\beta))=-\tan(\alpha+\beta)=-\frac{\tan\alpha+\tan\beta}{1-\tan\alpha\tan\beta}$.
  \item Multiplicando por el denominador: $\tan\gamma-\tan\alpha\tan\beta\tan\gamma=-\tan\alpha-\tan\beta$.
\end{enumerate}
\resultado $\mathbf{\tan\alpha+\tan\beta+\tan\gamma=\tan\alpha\tan\beta\tan\gamma}$.
\comprobacion triángulo equilátero: $3\sqrt3=(\sqrt3)^{3}$.
\end{solucion}

\begin{solucion}{eje:2-chebyshev}
\datos \cref{cor:suma}; $n\ge1$.
\planteamiento probar la recurrencia y aplicar inducción fuerte con los casos base $n=1,2$.
\pasos
\begin{enumerate}[label=\arabic*.]
  \item $\cos(n+1)x+\cos(n-1)x=(\cos nx\cos x-\sin nx\sin x)+(\cos nx\cos x+\sin nx\sin x)=2\cos x\cos nx$.
  \item Base: $T_1(t)=t$ (grado~1, coeficiente $1=2^{0}$) y $T_2(t)=2t^{2}-1$ (grado~2, coeficiente $2=2^{1}$), por la \cref{prop:doble}.
  \item Paso: si $T_n$ y $T_{n-1}$ existen con las propiedades indicadas, se define $T_{n+1}(t)=2tT_n(t)-T_{n-1}(t)$. Por el paso~1, $\cos(n+1)x=T_{n+1}(\cos x)$.
  \item $2tT_n(t)$ tiene grado $n+1$ y coeficiente principal $2\cdot2^{n-1}=2^{n}$; restar $T_{n-1}$, de grado $n-1$, no los altera.
\end{enumerate}
\resultado \textbf{$\cos nx=T_n(\cos x)$ con $T_n$ de grado $n$ y coeficiente principal $2^{n-1}$} (polinomios de Chebyshev).
\comprobacion $T_3=2t(2t^{2}-1)-t=4t^{3}-3t$, que coincide con el \cref{eje:2-triple}.
\end{solucion}

\begin{solucion}{eje:2-potencia-giro}
\datos $R(\alpha)R(\beta)=R(\alpha+\beta)$ (\cref{cor:giros}).
\planteamiento inducción para $n\ge0$; inversa a partir de $R(0)=I$; potencias negativas como inversas.
\pasos
\begin{enumerate}[label=\arabic*.]
  \item $n=0$: $R(\alpha)^{0}=I=R(0)$, porque $\cos0=1$ y $\sin0=0$.
  \item Si $R(\alpha)^{n}=R(n\alpha)$, entonces $R(\alpha)^{n+1}=R(n\alpha)R(\alpha)=R((n+1)\alpha)$.
  \item $R(\alpha)R(-\alpha)=R(0)=I=R(-\alpha)R(\alpha)$, luego $R(\alpha)^{-1}=R(-\alpha)$. Por la paridad, $R(-\alpha)=\begin{pmatrix}\cos\alpha&\sin\alpha\\-\sin\alpha&\cos\alpha\end{pmatrix}=R(\alpha)\T$.
  \item Para $n=-m<0$: $R(\alpha)^{-m}=(R(\alpha)^{m})^{-1}=R(m\alpha)^{-1}=R(-m\alpha)$.
\end{enumerate}
\resultado \textbf{$R(\alpha)^{n}=R(n\alpha)$ para todo $n\in\Z$, y $R(\alpha)^{-1}=R(-\alpha)=R(\alpha)\T$.}
\comprobacion $R(\pi/2)^{4}=R(2\pi)=I$: cuatro giros rectos dan una vuelta completa.
\end{solucion}

% ---------------------------------------------------------------------
\section{Tema 3}

\begin{solucion}{eje:3-que-teorema}
\planteamiento identificar el caso de la \cref{tab:casos}.
\pasos
\begin{enumerate}[label=(\alph*)]
  \item LLL: teorema del coseno; $4+6=10>9$, luego hay \textbf{un} triángulo.
  \item ALA ($c$ es el lado comprendido entre $A$ y $B$): $C=65^{\circ}$ y teorema del seno; \textbf{uno}.
  \item LAL: teorema del coseno para $a$; \textbf{uno}.
  \item LLA: teorema del seno; $h=8\sin30^{\circ}=4<5<8$, luego \textbf{dos} (\cref{prop:ambiguo}).
\end{enumerate}
\resultado \textbf{(a) coseno, una; (b) seno, una; (c) coseno, una; (d) seno, dos.}
\end{solucion}

\begin{solucion}{eje:3-imposible}
\planteamiento desigualdad triangular para el primero; \cref{prop:ambiguo} para el segundo.
\pasos
\begin{enumerate}[label=\arabic*.]
  \item $3+4=7<8$: el lado mayor supera la suma de los otros dos; no existe el triángulo.
  \item $h=b\sin A=5\cdot\frac{\sqrt3}{2}\approx4{,}330>a=2$: no existe. Equivalentemente, $\sin B=\frac{5\sin60^{\circ}}{2}\approx2{,}165>1$.
\end{enumerate}
\resultado \textbf{Ninguno de los dos triángulos existe.}
\end{solucion}

\begin{solucion}{eje:3-obtuso}
\datos lados $6$, $7$, $10$.
\planteamiento \cref{cor:pitagoras} con el lado mayor.
\pasos
\begin{enumerate}[label=\arabic*.]
  \item $10^{2}=100$ y $6^{2}+7^{2}=85$.
  \item $100>85$: el ángulo opuesto al lado~10 es obtuso.
\end{enumerate}
\resultado \textbf{Obtusángulo.}
\comprobacion $\cos C=\frac{85-100}{84}=-\frac{15}{84}$, y $C\approx100{,}3^{\circ}$.
\end{solucion}

\begin{solucion}{eje:3-ala}
\datos $A=45^{\circ}$, $B=60^{\circ}$, $c=10$.
\planteamiento $C=180^{\circ}-A-B$; teorema del seno; área con la \cref{prop:area}.
\pasos
\begin{enumerate}[label=\arabic*.]
  \item $C=75^{\circ}$, con $\sin75^{\circ}=\frac{\sqrt6+\sqrt2}{4}$ (\cref{eje:2-exactos}).
  \item $a=\frac{c\sin A}{\sin C}=\frac{10\cdot\sqrt2/2}{(\sqrt6+\sqrt2)/4}=\frac{20}{\sqrt3+1}=10(\sqrt3-1)\approx7{,}321$.
  \item $b=\frac{c\sin B}{\sin C}=\frac{20\sqrt3}{\sqrt6+\sqrt2}=15\sqrt2-5\sqrt6\approx8{,}966$.
  \item $S=\frac12\,c\,a\sin B=\frac12\cdot10\cdot10(\sqrt3-1)\cdot\frac{\sqrt3}{2}=25(3-\sqrt3)\approx31{,}70$.
\end{enumerate}
\resultado \textbf{$C=75^{\circ}$, $a=10(\sqrt3-1)\approx7{,}321$, $b\approx8{,}966$, $S=75-25\sqrt3\approx31{,}70$.}
\comprobacion $a/\sin A=7{,}321/0{,}7071=10{,}35$; $b/\sin B=8{,}966/0{,}8660=10{,}35$; $c/\sin C=10/0{,}9659=10{,}35$.
\end{solucion}

\begin{solucion}{eje:3-ambiguo}
\datos $a=6$, $b=8$, $A=35^{\circ}$, $\sin35^{\circ}\approx0{,}5736$.
\planteamiento \cref{prop:ambiguo} y teorema del seno para cada solución.
\pasos
\begin{enumerate}[label=\arabic*.]
  \item $h=8\cdot0{,}5736=4{,}589$; como $h<6<8$, hay \textbf{dos} triángulos.
  \item $\sin B=\frac{8\cdot0{,}5736}{6}=0{,}7648$: $B_1=\arcsin0{,}7648\approx49{,}89^{\circ}$ y $B_2=180^{\circ}-B_1\approx130{,}11^{\circ}$.
  \item Primer triángulo: $C_1=180^{\circ}-35^{\circ}-49{,}89^{\circ}=95{,}11^{\circ}$; $c_1=\frac{6\sin C_1}{\sin35^{\circ}}\approx10{,}42$.
  \item Segundo: $C_2=180^{\circ}-35^{\circ}-130{,}11^{\circ}=14{,}89^{\circ}$; $c_2=\frac{6\sin C_2}{\sin35^{\circ}}\approx2{,}687$.
\end{enumerate}
\resultado \textbf{$(B,C,c)\approx(49{,}89^{\circ},\,95{,}11^{\circ},\,10{,}42)$ y $(130{,}11^{\circ},\,14{,}89^{\circ},\,2{,}687)$.}
\comprobacion teorema del coseno en el segundo: $2{,}687^{2}+8^{2}-2\cdot2{,}687\cdot8\cos35^{\circ}\approx7{,}22+64-35{,}22=36{,}0=a^{2}$.
\end{solucion}

\begin{solucion}{eje:3-lal}
\datos $b=5$, $c=7$, $A=60^{\circ}$.
\planteamiento teorema del coseno para $a$ y para un ángulo; el tercero, por diferencia.
\pasos
\begin{enumerate}[label=\arabic*.]
  \item $a^{2}=25+49-2\cdot5\cdot7\cdot\frac12=39$; $a=\sqrt{39}\approx6{,}245$.
  \item $\cos B=\frac{a^{2}+c^{2}-b^{2}}{2ac}=\frac{39+49-25}{14\sqrt{39}}=\frac{63}{14\sqrt{39}}\approx0{,}7206$; $B\approx43{,}90^{\circ}$.
  \item $C=180^{\circ}-60^{\circ}-43{,}90^{\circ}=76{,}10^{\circ}$.
\end{enumerate}
\resultado \textbf{$a=\sqrt{39}\approx6{,}245$, $B\approx43{,}90^{\circ}$, $C\approx76{,}10^{\circ}$.}
\comprobacion teorema del seno: $a/\sin A=6{,}245/0{,}8660=7{,}211$ y $c/\sin C=7/0{,}9707=7{,}211$.
\end{solucion}

\begin{solucion}{eje:3-barcos}
\datos lados $3$ y $5$ km desde el puerto; ángulo comprendido $120^{\circ}$.
\planteamiento caso LAL: teorema del coseno para la distancia y después para el ángulo pedido.
\pasos
\begin{enumerate}[label=\arabic*.]
  \item $d^{2}=3^{2}+5^{2}-2\cdot3\cdot5\cos120^{\circ}=9+25+15=49$; $d=7$ km.
  \item En el segundo barco concurren los lados $5$ (al puerto) y $7$ (al primer barco), y el opuesto mide~3: $\cos\varphi=\frac{25+49-9}{2\cdot5\cdot7}=\frac{65}{70}=\frac{13}{14}$; $\varphi\approx21{,}79^{\circ}$.
\end{enumerate}
\resultado \textbf{$7$ km; ángulo $\approx21{,}79^{\circ}$.}
\comprobacion el ángulo en el primer barco es $\arccos\frac{9+49-25}{42}=\arccos\frac{33}{42}\approx38{,}21^{\circ}$, y $120^{\circ}+21{,}79^{\circ}+38{,}21^{\circ}=180^{\circ}$.
\end{solucion}

\begin{solucion}{eje:3-lll}
\datos $a=4$, $b=5$, $c=6$.
\planteamiento teorema del coseno para cada ángulo; $\cos2A$ con la \cref{prop:doble}; área por los dos métodos.
\pasos
\begin{enumerate}[label=\arabic*.]
  \item $\cos A=\frac{25+36-16}{60}=\frac34$, $A\approx41{,}41^{\circ}$; $\cos B=\frac{16+36-25}{48}=\frac{9}{16}$, $B\approx55{,}77^{\circ}$; $\cos C=\frac{16+25-36}{40}=\frac18$, $C\approx82{,}82^{\circ}$.
  \item $\cos2A=2\cdot\frac{9}{16}-1=\frac18=\cos C$. Como $2A$ y $C$ están en $(0,\pi)$, donde el coseno es inyectivo, $C=2A$.
  \item $\sin C=\sqrt{1-\frac1{64}}=\frac{3\sqrt7}{8}$; $S=\frac12\cdot4\cdot5\cdot\frac{3\sqrt7}{8}=\frac{15\sqrt7}{4}$.
  \item Herón: $s=\frac{15}{2}$; $s(s-a)(s-b)(s-c)=\frac{15}{2}\cdot\frac72\cdot\frac52\cdot\frac32=\frac{1575}{16}$; $S=\frac{\sqrt{1575}}{4}=\frac{15\sqrt7}{4}$.
\end{enumerate}
\resultado \textbf{$A\approx41{,}41^{\circ}$, $B\approx55{,}77^{\circ}$, $C=2A\approx82{,}82^{\circ}$; $S=\frac{15\sqrt7}{4}\approx9{,}922$.}
\comprobacion $41{,}41+55{,}77+82{,}82=180{,}00$ y $1575=225\cdot7$.
\end{solucion}

\begin{solucion}{eje:3-angulo-vectores}
\datos $u=(1,2,2)$, $v=(2,-1,2)$.
\planteamiento \cref{teo:escalar} y \cref{cor:escalar}(iii).
\pasos
\begin{enumerate}[label=\arabic*.]
  \item $u\cdot v=2-2+4=4$; $\norm{u}=\sqrt{9}=3$; $\norm{v}=\sqrt9=3$.
  \item $\cos\varphi=\frac49$; $\varphi\approx63{,}61^{\circ}$.
  \item Componente: $u\cdot v/\norm{v}=\frac43$; proyección: $\frac{4}{9}(2,-1,2)=(\frac89,-\frac49,\frac89)$.
\end{enumerate}
\resultado \textbf{$\varphi=\arccos\frac49\approx63{,}61^{\circ}$; componente $\frac43$; proyección $(\frac89,-\frac49,\frac89)$.}
\comprobacion $u-\mathrm{proy}=(\frac19,\frac{22}{9},\frac{10}{9})$ es ortogonal a $v$: $\frac29-\frac{22}{9}+\frac{20}{9}=0$.
\end{solucion}

\begin{solucion}{eje:3-coseno-normas}
\datos $\norm{u}=3$, $\norm{v}=4$, $\norm{u-v}=\sqrt{13}$.
\planteamiento ecuación~\eqref{eq:norma-diferencia} y su análoga para $u+v$.
\pasos
\begin{enumerate}[label=\arabic*.]
  \item $13=9+16-2u\cdot v$, luego $u\cdot v=6$.
  \item $\cos\varphi=\frac{6}{12}=\frac12$, $\varphi=\pi/3$.
  \item $\norm{u+v}^{2}=9+16+2\cdot6=37$, $\norm{u+v}=\sqrt{37}\approx6{,}083$.
\end{enumerate}
\resultado \textbf{$u\cdot v=6$, $\varphi=60^{\circ}$, $\norm{u+v}=\sqrt{37}$.}
\comprobacion ley del paralelogramo: $\norm{u+v}^{2}+\norm{u-v}^{2}=37+13=50=2(9+16)$.
\end{solucion}

\begin{solucion}{eje:3-similitud}
\datos $a=(1,0,1,1)$, $b=(0,1,1,1)$, $c=(2,0,2,2)$.
\planteamiento $\cos(x,y)=x\cdot y/(\norm{x}\norm{y})$.
\pasos
\begin{enumerate}[label=\arabic*.]
  \item $\norm{a}=\norm{b}=\sqrt3$, $\norm{c}=\sqrt{12}=2\sqrt3$.
  \item $a\cdot b=2$: $\cos(a,b)=\frac{2}{3}$, ángulo $\approx48{,}19^{\circ}$.
  \item $a\cdot c=6$: $\cos(a,c)=\frac{6}{\sqrt3\cdot2\sqrt3}=1$, ángulo $0$.
  \item $b\cdot c=4$: $\cos(b,c)=\frac{4}{6}=\frac23$, ángulo $\approx48{,}19^{\circ}$.
\end{enumerate}
\resultado \textbf{$\frac23$, $1$, $\frac23$.} Como $c=2a$, ambos vectores tienen la misma dirección y la similitud coseno, que ignora la longitud, vale~1.
\comprobacion $\cos(b,c)=\cos(b,2a)=\cos(b,a)$, como debe ser.
\end{solucion}

\begin{solucion}{eje:3-circunradio}
\datos triángulo $ABC$ inscrito en una circunferencia de centro $O$ y radio $R$.
\planteamiento construir un triángulo rectángulo con hipotenusa $2R$ y cateto $a$.
\pasos
\begin{enumerate}[label=\arabic*.]
  \item Sea $D$ el punto diametralmente opuesto a $B$, de modo que $BD=2R$. Si $D=C$, $BC$ es un diámetro, el ángulo $A$ es recto (está inscrito en una semicircunferencia) y $a=2R=2R\sin A$.
  \item Si $D\ne C$, el triángulo $BCD$ es rectángulo en $C$ (ángulo inscrito en una semicircunferencia), luego $a=BC=BD\sin\angle BDC=2R\sin\angle BDC$.
  \item Si $A$ es agudo, $A$ y $D$ están en el mismo arco determinado por $BC$, y el teorema del ángulo inscrito da $\angle BDC=A$.
  \item Si $A$ es obtuso, $A$ y $D$ están en arcos opuestos; el cuadrilátero $ABDC$ es inscriptible y sus ángulos opuestos suman $\pi$: $\angle BDC=\pi-A$, con $\sin(\pi-A)=\sin A$.
\end{enumerate}
\resultado \textbf{$a=2R\sin A$ en todos los casos}; junto con el \cref{teo:seno}, $\frac{a}{\sin A}=\frac{b}{\sin B}=\frac{c}{\sin C}=2R$.
\comprobacion triángulo equilátero de lado~1: $2R=1/\sin60^{\circ}=2/\sqrt3$, y en efecto el circunradio es $1/\sqrt3$.
\end{solucion}

\begin{solucion}{eje:3-heron}
\datos $S=\frac12ab\sin C$; $c^{2}=a^{2}+b^{2}-2ab\cos C$; $s=\frac{a+b+c}{2}$.
\planteamiento eliminar el ángulo y factorizar dos veces como diferencia de cuadrados.
\pasos
\begin{enumerate}[label=\arabic*.]
  \item $16S^{2}=4a^{2}b^{2}\sin^{2}C=(2ab)^{2}-(2ab\cos C)^{2}=(2ab)^{2}-(a^{2}+b^{2}-c^{2})^{2}$.
  \item $=\bigl(2ab-a^{2}-b^{2}+c^{2}\bigr)\bigl(2ab+a^{2}+b^{2}-c^{2}\bigr)=\bigl(c^{2}-(a-b)^{2}\bigr)\bigl((a+b)^{2}-c^{2}\bigr)$.
  \item $=(c-a+b)(c+a-b)(a+b-c)(a+b+c)$.
  \item Con $a+b+c=2s$: $c-a+b=2(s-a)$, $c+a-b=2(s-b)$, $a+b-c=2(s-c)$; el producto es $16\,s(s-a)(s-b)(s-c)$.
\end{enumerate}
\resultado $\mathbf{S=\sqrt{s(s-a)(s-b)(s-c)}}$.
\comprobacion triángulo $3$, $4$, $5$: $s=6$, $\sqrt{6\cdot3\cdot2\cdot1}=6=\frac12\cdot3\cdot4$.
\end{solucion}

\begin{solucion}{eje:3-proyecciones}
\datos coordenadas de la demostración del \cref{teo:coseno}: $C=(0,0)$, $B=(a,0)$, $A=(b\cos C,b\sin C)$.
\planteamiento calcular la abscisa de $A$ de dos maneras.
\pasos
\begin{enumerate}[label=\arabic*.]
  \item Desde $C$: $x_A=b\cos C$.
  \item Desde $B$: el lado $BA$, de longitud $c$, forma con el semieje negativo de abscisas el ángulo $B$, luego $A=B+c(-\cos B,\sin B)$ y $x_A=a-c\cos B$.
  \item Igualando: $a=b\cos C+c\cos B$, en todos los casos (los cosenos negativos de los ángulos obtusos dan el signo correcto).
  \item (b) Las tres leyes son $a=b\cos C+c\cos B$, $b=c\cos A+a\cos C$ y $c=a\cos B+b\cos A$. Multiplicando por $a$, $b$ y $c$: $a^{2}=ab\cos C+ac\cos B$, $b^{2}=bc\cos A+ab\cos C$, $c^{2}=ac\cos B+bc\cos A$.
  \item $b^{2}+c^{2}-a^{2}=2bc\cos A$.
\end{enumerate}
\resultado \textbf{$a=b\cos C+c\cos B$, y de las tres leyes se obtiene $a^{2}=b^{2}+c^{2}-2bc\cos A$.}
\comprobacion triángulo $3,4,5$ rectángulo en $C$ ($c=5$): $c=a\cos B+b\cos A=3\cdot\frac35+4\cdot\frac45=5$.
\end{solucion}

\begin{solucion}{eje:3-mayor-lado}
\datos $A+B<\pi$; teorema del seno.
\planteamiento $a>b\iff\sin A>\sin B$, y el signo de $\sin A-\sin B$ se lee en su forma de producto.
\pasos
\begin{enumerate}[label=\arabic*.]
  \item Por el \cref{teo:seno}, $a/b=\sin A/\sin B$, luego $a>b\iff\sin A>\sin B$.
  \item $\sin A-\sin B=2\cos\frac{A+B}{2}\sin\frac{A-B}{2}$.
  \item $0<A+B<\pi$, luego $\frac{A+B}{2}\in(0,\frac{\pi}{2})$ y el coseno es positivo.
  \item $\frac{A-B}{2}\in(-\frac{\pi}{2},\frac{\pi}{2})$, donde el seno tiene el signo de su argumento.
  \item Por tanto, $\sin A-\sin B>0\iff A-B>0$.
\end{enumerate}
\resultado \textbf{$a>b\iff A>B$.}
\comprobacion en el \cref{eje:3-lll}, $4<5<6$ y $41{,}41^{\circ}<55{,}77^{\circ}<82{,}82^{\circ}$.
\end{solucion}

% ---------------------------------------------------------------------
\section{Tema 4}

\begin{solucion}{eje:4-parametros}
\datos $y=3\sin(2x-\frac{\pi}{3})+1$.
\planteamiento escribir $2x-\frac{\pi}{3}=2(x-\frac{\pi}{6})$ y leer los parámetros de la sinusoide.
\pasos
\begin{enumerate}[label=\arabic*.]
  \item $A=3$, $\omega=2$, $x_0=\pi/6$, $D=1$.
  \item Periodo $2\pi/2=\pi$; amplitud $3$.
  \item Máximo $1+3=4$, mínimo $1-3=-2$.
\end{enumerate}
\resultado \textbf{amplitud 3, periodo $\pi$, desplazamiento horizontal $\pi/6$ a la derecha, desplazamiento vertical 1, máximo 4 y mínimo $-2$.}
\comprobacion el máximo se alcanza donde $2x-\frac{\pi}{3}=\frac{\pi}{2}$, es decir, en $x=\frac{5\pi}{12}$: $3\sin\frac{\pi}{2}+1=4$.
\end{solucion}

\begin{solucion}{eje:4-dominios}
\planteamiento \cref{def:inversas} y \cref{prop:inversas}(i).
\pasos
\begin{enumerate}[label=\arabic*.]
  \item $\arcsin\colon[-1,1]\to[-\frac{\pi}{2},\frac{\pi}{2}]$; $\arccos\colon[-1,1]\to[0,\pi]$; $\arctan\colon\R\to(-\frac{\pi}{2},\frac{\pi}{2})$.
  \item El seno no es inyectivo en $\R$ (por ejemplo, $\sin0=\sin\pi$), y solo las funciones inyectivas tienen inversa; se elige un intervalo donde lo es y donde alcanza todos sus valores.
  \item $\arcsin(\sin\frac{3\pi}{4})=\arcsin\frac{\sqrt2}{2}=\frac{\pi}{4}$.
  \item $\arccos(\cos(-\frac{\pi}{3}))=\arccos\frac12=\frac{\pi}{3}$.
  \item $\tan(\arctan5)=5$, porque $\arctan$ es la inversa de la tangente restringida.
\end{enumerate}
\resultado \textbf{$\frac{\pi}{4}$, $\frac{\pi}{3}$ y $5$.}
\comprobacion $\frac{\pi}{4}\in[-\frac{\pi}{2},\frac{\pi}{2}]$ tiene el mismo seno que $\frac{3\pi}{4}$, su suplementario.
\end{solucion}

\begin{solucion}{eje:4-exactos}
\planteamiento \cref{tab:notables}, rangos de la \cref{def:inversas} y \cref{prop:inversas}(ii).
\pasos
\begin{enumerate}[label=\arabic*.]
  \item $\arcsin(-\frac12)=-\frac{\pi}{6}$ (impar).
  \item $\arccos(-\frac{\sqrt2}{2})=\pi-\frac{\pi}{4}=\frac{3\pi}{4}$.
  \item $\arctan\sqrt3=\frac{\pi}{3}$.
  \item $\cos\frac{5\pi}{3}=\frac12$, luego $\arccos(\cos\frac{5\pi}{3})=\frac{\pi}{3}$.
  \item $\sin(\arccos\frac35)=\sqrt{1-\frac{9}{25}}=\frac45$.
  \item Sea $\theta=\arcsin\frac13$: $\cos\theta=\sqrt{1-\frac19}=\frac{2\sqrt2}{3}$ y $\tan\theta=\frac{1/3}{2\sqrt2/3}=\frac{1}{2\sqrt2}=\frac{\sqrt2}{4}$.
\end{enumerate}
\resultado $\mathbf{-\frac{\pi}{6},\ \frac{3\pi}{4},\ \frac{\pi}{3},\ \frac{\pi}{3},\ \frac45,\ \frac{\sqrt2}{4}}$.
\comprobacion $\tan(\arcsin\frac13)=\tan(0{,}3398)=0{,}3536=\sqrt2/4$.
\end{solucion}

\begin{solucion}{eje:4-basicas}
\planteamiento \cref{teo:ecuaciones}.
\pasos
\begin{enumerate}[label=\arabic*.]
  \item (a) $\sin x=\frac12$, $\alpha=\frac{\pi}{6}$: $x=\frac{\pi}{6}+2k\pi$ o $x=\frac{5\pi}{6}+2k\pi$. En $[0,2\pi)$: $\frac{\pi}{6}$ y $\frac{5\pi}{6}$.
  \item (b) $\arctan(-\sqrt3)=-\frac{\pi}{3}$: $x=-\frac{\pi}{3}+k\pi$.
  \item (c) $2x=\pm\frac{\pi}{3}+2k\pi$, luego $x=\pm\frac{\pi}{6}+k\pi$. Con $x\in[0,2\pi)$: $\frac{\pi}{6}$, $\frac{5\pi}{6}$, $\frac{7\pi}{6}$, $\frac{11\pi}{6}$.
\end{enumerate}
\resultado \textbf{(a) $\{\frac{\pi}{6},\frac{5\pi}{6}\}$ y sus trasladados por $2k\pi$; (b) $-\frac{\pi}{3}+k\pi$; (c) $\{\frac{\pi}{6},\frac{5\pi}{6},\frac{7\pi}{6},\frac{11\pi}{6}\}$.}
\comprobacion en (c), $\cos\frac{5\pi}{3}=\frac12$ con $x=\frac{5\pi}{6}$, y $\cos\frac{7\pi}{3}=\cos\frac{\pi}{3}=\frac12$ con $x=\frac{7\pi}{6}$.
\end{solucion}

\begin{solucion}{eje:4-cuadratica}
\planteamiento sustituir $c=\cos x$ y factorizar.
\pasos
\begin{enumerate}[label=\arabic*.]
  \item $2c^{2}-c-1=(2c+1)(c-1)=0$: $c=1$ o $c=-\frac12$.
  \item $\cos x=1$: $x=2k\pi$.
  \item $\cos x=-\frac12$: $\arccos(-\frac12)=\frac{2\pi}{3}$, luego $x=\pm\frac{2\pi}{3}+2k\pi$.
\end{enumerate}
\resultado $\mathbf{x=2k\pi}$ \textbf{o} $\mathbf{x=\pm\frac{2\pi}{3}+2k\pi}$, $k\in\Z$.
\comprobacion $x=\frac{2\pi}{3}$: $2\cdot\frac14+\frac12-1=0$.
\end{solucion}

\begin{solucion}{eje:4-doble}
\planteamiento $\sin2x=2\sin x\cos x$; factorizar en lugar de dividir.
\pasos
\begin{enumerate}[label=\arabic*.]
  \item $2\sin x\cos x-\cos x=\cos x(2\sin x-1)=0$.
  \item $\cos x=0$: $x=\frac{\pi}{2},\frac{3\pi}{2}$.
  \item $\sin x=\frac12$: $x=\frac{\pi}{6},\frac{5\pi}{6}$.
\end{enumerate}
\resultado $\mathbf{x\in\{\frac{\pi}{6},\frac{\pi}{2},\frac{5\pi}{6},\frac{3\pi}{2}\}}$. Dividir por $\cos x$ pierde $\frac{\pi}{2}$ y $\frac{3\pi}{2}$.
\comprobacion $x=\frac{3\pi}{2}$: $\sin3\pi=0=\cos\frac{3\pi}{2}$.
\end{solucion}

\begin{solucion}{eje:4-armonico}
\datos $a=\sqrt3$, $b=1$.
\planteamiento \cref{prop:armonico} y \cref{teo:ecuaciones}.
\pasos
\begin{enumerate}[label=\arabic*.]
  \item $R=\sqrt{3+1}=2$; $\cos\varphi=\frac{\sqrt3}{2}$, $\sin\varphi=\frac12$, luego $\varphi=\frac{\pi}{6}$: $\sqrt3\sin x+\cos x=2\sin(x+\frac{\pi}{6})$.
  \item Máximo $2$ cuando $x+\frac{\pi}{6}=\frac{\pi}{2}$, es decir, $x=\frac{\pi}{3}$.
  \item $2\sin(x+\frac{\pi}{6})=1$: $x+\frac{\pi}{6}=\frac{\pi}{6}+2k\pi$ o $\frac{5\pi}{6}+2k\pi$, de donde $x=2k\pi$ o $x=\frac{2\pi}{3}+2k\pi$.
\end{enumerate}
\resultado \textbf{$2\sin(x+\frac{\pi}{6})$; máximo 2 en $x=\frac{\pi}{3}$; soluciones en $[0,2\pi)$: $0$ y $\frac{2\pi}{3}$.}
\comprobacion $x=\frac{2\pi}{3}$: $\sqrt3\cdot\frac{\sqrt3}{2}-\frac12=1$; $x=\frac{\pi}{3}$: $\frac32+\frac12=2$.
\end{solucion}

\begin{solucion}{eje:4-cuadrado}
\planteamiento (a) elevar al cuadrado y comprobar; (b) $\sin x+\cos x=\sqrt2\sin(x+\frac{\pi}{4})$.
\pasos
\begin{enumerate}[label=\arabic*.]
  \item (a) $(\sin x+\cos x)^{2}=1+\sin2x=1$, luego $\sin2x=0$ y $x\in\{0,\frac{\pi}{2},\pi,\frac{3\pi}{2}\}$.
  \item Sustituyendo en la ecuación original: $0\mapsto0+1=1$ sí; $\frac{\pi}{2}\mapsto1+0=1$ sí; $\pi\mapsto0-1=-1$ no; $\frac{3\pi}{2}\mapsto-1+0=-1$ no.
  \item (b) $\sin(x+\frac{\pi}{4})=\frac{\sqrt2}{2}$: $x+\frac{\pi}{4}\in\{\frac{\pi}{4},\frac{3\pi}{4}\}+2k\pi$, luego $x\in\{0,\frac{\pi}{2}\}$ en $[0,2\pi)$.
  \item Elevar al cuadrado no es reversible: $A^{2}=B^{2}$ equivale a $A=\pm B$, y las soluciones falsas son las de $\sin x+\cos x=-1$.
\end{enumerate}
\resultado $\mathbf{x\in\{0,\frac{\pi}{2}\}}$; $\pi$ y $\frac{3\pi}{2}$ son soluciones falsas introducidas al elevar al cuadrado.
\comprobacion ambos métodos coinciden tras descartar las soluciones falsas.
\end{solucion}

\begin{solucion}{eje:4-periodo-suma}
\datos $f(x)=\sin3x+\cos2x$; $g(x)=\sin^{2}x$.
\planteamiento para $f$: comprobar $2\pi$ y descartar periodos menores evaluando en $0$ y $\pm T$. Para $g$: reducción de potencias.
\pasos
\begin{enumerate}[label=\arabic*.]
  \item $f(x+2\pi)=\sin(3x+6\pi)+\cos(2x+4\pi)=f(x)$.
  \item Si $T>0$ es periodo, $f(T)=f(-T)=f(0)=1$: $\sin3T+\cos2T=1$ y $-\sin3T+\cos2T=1$. Sumando y restando, $\cos2T=1$ y $\sin3T=0$, luego $T\in\pi\Z$ (\cref{teo:simetrias}(ii)) y $T\in\frac{\pi}{3}\Z$ (\cref{eje:1-periodo-tan}(a)). Así $T\in\{\pi,2\pi,\dots\}$.
  \item $T=\pi$ no es periodo: $f(x+\pi)=-\sin3x+\cos2x$, y en $x=\frac{\pi}{6}$ vale $-1+\frac12=-\frac12\ne f(\frac{\pi}{6})=1+\frac12=\frac32$.
  \item $g(x)=\frac{1-\cos2x}{2}$; por la \cref{prop:periodos}(ii), $\cos2x$ tiene periodo fundamental $\pi$, y sumar una constante y multiplicar por $-\frac12$ no lo cambia.
\end{enumerate}
\resultado \textbf{$f$ tiene periodo fundamental $2\pi$; $g$, periodo fundamental $\pi$.}
\comprobacion $g(x+\pi)=(\sin(x+\pi))^{2}=(-\sin x)^{2}=g(x)$.
\end{solucion}

\begin{solucion}{eje:4-pe}
\datos $d=8$; $\omega_i=10000^{-i/4}$, $i=0,1,2,3$.
\planteamiento evaluar $\PE(p)_{2i}=\sin p\omega_i$ y $\PE(p)_{2i+1}=\cos p\omega_i$ en radianes.
\pasos
\begin{enumerate}[label=\arabic*.]
  \item $\omega_0=1$, $\omega_1=0{,}1$, $\omega_2=0{,}01$, $\omega_3=0{,}001$.
  \item Periodos $2\pi/\omega_i$: $6{,}283$; $62{,}83$; $628{,}3$; $6283$.
  \item $\PE(1)=(0{,}8415;\ 0{,}5403;\ 0{,}0998;\ 0{,}9950;\ 0{,}0100;\ 1{,}0000;\ 0{,}0010;\ 1{,}0000)$.
  \item $\PE(2)=(0{,}9093;\ -0{,}4161;\ 0{,}1987;\ 0{,}9801;\ 0{,}0200;\ 0{,}9998;\ 0{,}0020;\ 1{,}0000)$.
\end{enumerate}
\resultado \textbf{Frecuencias $1;\ 0{,}1;\ 0{,}01;\ 0{,}001$ y periodos $2\pi\cdot10^{i}$.} El primer par (periodo $6{,}28$) distingue claramente las posiciones $1$ y $2$; los pares lentos apenas cambian entre posiciones vecinas y sirven para distinguir posiciones lejanas.
\comprobacion cada par cumple $\sin^{2}+\cos^{2}=1$: $0{,}8415^{2}+0{,}5403^{2}=0{,}7081+0{,}2919=1{,}0000$.
\end{solucion}

\begin{solucion}{eje:4-creciente}
\datos $-\frac{\pi}{2}\le b<a\le\frac{\pi}{2}$.
\planteamiento $\sin a-\sin b=2\cos\frac{a+b}{2}\sin\frac{a-b}{2}$ (\cref{cor:suma-producto}).
\pasos
\begin{enumerate}[label=\arabic*.]
  \item $-\pi<a+b<\pi$ (no pueden ser ambos iguales a $\frac{\pi}{2}$ ni ambos a $-\frac{\pi}{2}$), luego $\frac{a+b}{2}\in(-\frac{\pi}{2},\frac{\pi}{2})$ y $\cos\frac{a+b}{2}>0$.
  \item $0<a-b\le\pi$, luego $\frac{a-b}{2}\in(0,\frac{\pi}{2}]$ y $\sin\frac{a-b}{2}>0$.
  \item El producto es positivo: $\sin a>\sin b$.
\end{enumerate}
\resultado \textbf{El seno es estrictamente creciente, y por tanto inyectivo, en $[-\frac{\pi}{2},\frac{\pi}{2}]$.}
\comprobacion $\sin(-\frac{\pi}{6})=-\frac12<\sin\frac{\pi}{6}=\frac12<\sin\frac{\pi}{2}=1$.
\end{solucion}

\begin{solucion}{eje:4-cos-iguales}
\planteamiento $\cos x-\cos y=-2\sin\frac{x+y}{2}\sin\frac{x-y}{2}$ y \cref{eje:1-periodo-tan}(a).
\pasos
\begin{enumerate}[label=\arabic*.]
  \item $\cos x=\cos y$ si y solo si $\sin\frac{x+y}{2}=0$ o $\sin\frac{x-y}{2}=0$.
  \item $\sin\frac{x+y}{2}=0\iff\frac{x+y}{2}=k\pi\iff x=-y+2k\pi$.
  \item $\sin\frac{x-y}{2}=0\iff\frac{x-y}{2}=k\pi\iff x=y+2k\pi$.
\end{enumerate}
\resultado $\mathbf{\cos x=\cos y\iff x=\pm y+2k\pi}$ para algún $k\in\Z$.
\comprobacion es el \cref{teo:ecuaciones}(ii) con $a=\cos y$, demostrado ahora por otra vía.
\end{solucion}

\begin{solucion}{eje:4-arctan-inverso}
\datos $x\ne0$.
\planteamiento para $x>0$, comprobar que $\frac{\pi}{2}-\arctan x$ está en el rango de $\arctan$ y tiene tangente $1/x$; para $x<0$, usar la imparidad.
\pasos
\begin{enumerate}[label=\arabic*.]
  \item $x>0$: $\theta=\arctan x\in(0,\frac{\pi}{2})$, luego $\frac{\pi}{2}-\theta\in(0,\frac{\pi}{2})$.
  \item $\tan(\frac{\pi}{2}-\theta)=\frac{\cos\theta}{\sin\theta}=\frac1{\tan\theta}=\frac1x$ (\cref{teo:simetrias}(vi)).
  \item Como $\frac{\pi}{2}-\theta\in(-\frac{\pi}{2},\frac{\pi}{2})$, $\arctan\frac1x=\frac{\pi}{2}-\theta$.
  \item $x<0$: $\arctan x+\arctan\frac1x=-\bigl(\arctan(-x)+\arctan\frac{1}{-x}\bigr)=-\frac{\pi}{2}$.
\end{enumerate}
\resultado $\mathbf{\arctan x+\arctan\frac1x=\frac{\pi}{2}\operatorname{signo}(x)}$.
\comprobacion $x=1$: $\frac{\pi}{4}+\frac{\pi}{4}=\frac{\pi}{2}$; $x=\sqrt3$: $\frac{\pi}{3}+\frac{\pi}{6}=\frac{\pi}{2}$.
\end{solucion}

\begin{solucion}{eje:4-machin}
\datos $\alpha=\arctan\frac12$, $\beta=\arctan\frac13$.
\planteamiento acotar $\alpha+\beta$ y calcular $\tan(\alpha+\beta)$.
\pasos
\begin{enumerate}[label=\arabic*.]
  \item La tangente es estrictamente creciente en $(-\frac{\pi}{2},\frac{\pi}{2})$ (en $[0,\frac{\pi}{2})$ el seno crece y el coseno, positivo, decrece; en el resto, por imparidad), luego también lo es $\arctan$. Como $0<\frac13<\frac12<1$: $0<\beta<\alpha<\arctan1=\frac{\pi}{4}$ y $\alpha+\beta\in(0,\frac{\pi}{2})$.
  \item $\tan(\alpha+\beta)=\frac{\frac12+\frac13}{1-\frac16}=\frac{5/6}{5/6}=1$.
  \item En $(-\frac{\pi}{2},\frac{\pi}{2})$ la tangente es inyectiva, y $\tan\frac{\pi}{4}=1$; luego $\alpha+\beta=\frac{\pi}{4}$.
  \item Sin la acotación, $\tan(\alpha+\beta)=1$ solo daría $\alpha+\beta=\frac{\pi}{4}+k\pi$ (\cref{teo:ecuaciones}(iii)).
\end{enumerate}
\resultado $\mathbf{\arctan\frac12+\arctan\frac13=\frac{\pi}{4}}$.
\comprobacion $0{,}463\,65+0{,}321\,75=0{,}785\,40\approx\pi/4$.
\end{solucion}

\begin{solucion}{eje:4-pe-lineal}
\datos $v(p)=(\sin\omega p,\cos\omega p)\T$.
\planteamiento desarrollar $\sin\omega(p+k)$ y $\cos\omega(p+k)$ con el \cref{cor:suma}.
\pasos
\begin{enumerate}[label=\arabic*.]
  \item $\sin\omega(p+k)=\cos\omega k\cdot\sin\omega p+\sin\omega k\cdot\cos\omega p$.
  \item $\cos\omega(p+k)=-\sin\omega k\cdot\sin\omega p+\cos\omega k\cdot\cos\omega p$.
  \item Luego $v(p+k)=M_kv(p)$ con $M_k=\begin{pmatrix}\cos\omega k&\sin\omega k\\-\sin\omega k&\cos\omega k\end{pmatrix}=R(-\omega k)$, independiente de $p$.
  \item $M_k\T M_k=R(\omega k)R(-\omega k)=R(0)=I$ (\cref{eje:2-potencia-giro}).
  \item $M_kM_l=R(-\omega k)R(-\omega l)=R(-\omega(k+l))=M_{k+l}$ (\cref{cor:giros}).
\end{enumerate}
\resultado \textbf{$M_k=R(-\omega k)$, ortogonal y con $M_kM_l=M_{k+l}$.}
\comprobacion $\omega=1$, $p=0$, $k=\frac{\pi}{2}$: $v(0)=(0,1)\T$ y $M_{\pi/2}v(0)=\begin{pmatrix}0&1\\-1&0\end{pmatrix}\begin{pmatrix}0\\1\end{pmatrix}=\begin{pmatrix}1\\0\end{pmatrix}=v(\frac{\pi}{2})$.
\end{solucion}

% ---------------------------------------------------------------------
\section{Tema 5}

\begin{solucion}{eje:5-partes}
\planteamiento aplicar las definiciones a cada número.
\pasos
\begin{enumerate}[label=\arabic*.]
  \item $z_1=3-4i$: $\Real=3$, $\Imag=-4$, $\bar z_1=3+4i$, $\abs{z_1}=\sqrt{9+16}=5$; cuadrante IV.
  \item $z_2=-2i$: $\Real=0$, $\Imag=-2$, $\bar z_2=2i$, $\abs{z_2}=2$; semieje imaginario negativo.
  \item $z_3=5$: $\Real=5$, $\Imag=0$, $\bar z_3=5$, $\abs{z_3}=5$; semieje real positivo.
  \item $z_4=-1+i$: $\Real=-1$, $\Imag=1$, $\bar z_4=-1-i$, $\abs{z_4}=\sqrt2$; cuadrante II.
\end{enumerate}
\resultado \textbf{módulos $5$, $2$, $5$ y $\sqrt2$}; las partes y los conjugados, como en el desarrollo.
\comprobacion $z_1\bar z_1=9+16=25=\abs{z_1}^{2}$.
\end{solucion}

\begin{solucion}{eje:5-vf}
\pasos
\begin{enumerate}[label=(\alph*)]
  \item \textbf{Falsa.} $z=1$, $w=-1$: $\abs{z+w}=0\ne2$.
  \item \textbf{Verdadera.} \Cref{prop:conjugado}(iv).
  \item \textbf{Verdadera.} \Cref{prop:conjugado}(i) con $w=z$.
  \item \textbf{Falsa.} La parte imaginaria es el real $2$.
\end{enumerate}
\resultado \textbf{(a) falsa; (b) verdadera; (c) verdadera; (d) falsa.}
\end{solucion}

\begin{solucion}{eje:5-potencias-i}
\planteamiento $i^{n}$ depende del resto de $n$ entre~4 (\cref{ej:potencias-i}).
\pasos
\begin{enumerate}[label=\arabic*.]
  \item $7=4+3$: $i^{7}=-i$. \quad $100=4\cdot25$: $i^{100}=1$.
  \item $i^{-1}=-i$, porque $i\cdot(-i)=-i^{2}=1$.
  \item $2026=4\cdot506+2$: $i^{2026}=-1$.
  \item $\sum_{k=0}^{10}i^{k}$: los ocho primeros términos forman dos bloques $1+i-1-i=0$; quedan $i^{8}+i^{9}+i^{10}=1+i-1=i$.
\end{enumerate}
\resultado $\mathbf{-i,\ 1,\ -i,\ -1}$ \textbf{y la suma vale} $\mathbf{i}$.
\comprobacion suma geométrica: $\frac{i^{11}-1}{i-1}=\frac{-i-1}{i-1}=\frac{(-1-i)(-1-i)}{(i-1)(-1-i)}=\frac{2i}{2}=i$.
\end{solucion}

\begin{solucion}{eje:5-operaciones}
\planteamiento operar con $i^{2}=-1$; para dividir, multiplicar por el conjugado del denominador.
\pasos
\begin{enumerate}[label=\arabic*.]
  \item $(3+2i)(1-4i)=3-12i+2i-8i^{2}=11-10i$.
  \item $(2-i)^{2}=4-4i+i^{2}=3-4i$.
  \item $\frac{2-i}{1+3i}=\frac{(2-i)(1-3i)}{1+9}=\frac{2-6i-i+3i^{2}}{10}=\frac{-1-7i}{10}$.
  \item $(1+i)^{2}=2i$, luego $(1+i)^{4}=(2i)^{2}=-4$.
\end{enumerate}
\resultado $\mathbf{11-10i,\ 3-4i,\ -\frac{1}{10}-\frac{7}{10}i,\ -4}$.
\comprobacion módulos: $\abs{3+2i}\abs{1-4i}=\sqrt{13}\sqrt{17}=\sqrt{221}$ y $\abs{11-10i}=\sqrt{121+100}=\sqrt{221}$.
\end{solucion}

\begin{solucion}{eje:5-inversos}
\planteamiento \cref{teo:cuerpo} y \cref{prop:modulo}(iii).
\pasos
\begin{enumerate}[label=\arabic*.]
  \item $i^{-1}=\frac{\bar i}{\abs{i}^{2}}=-i$.
  \item $(1+i)^{-1}=\frac{1-i}{2}$.
  \item $\frac{1+i}{2-i}=\frac{(1+i)(2+i)}{5}=\frac{2+i+2i+i^{2}}{5}=\frac{1+3i}{5}$.
  \item $\abs{z}\,\abs{z^{-1}}=\abs{zz^{-1}}=1$, luego $\abs{z^{-1}}=1/\abs{z}$.
\end{enumerate}
\resultado $\mathbf{-i,\ \frac{1-i}{2},\ \frac15+\frac35i}$ \textbf{y} $\mathbf{\abs{z^{-1}}=1/\abs{z}}$.
\comprobacion $(1+i)\cdot\frac{1-i}{2}=\frac{1-i^{2}}{2}=1$.
\end{solucion}

\begin{solucion}{eje:5-cuadraticas}
\planteamiento \cref{prop:cuadratica}.
\pasos
\begin{enumerate}[label=\arabic*.]
  \item $x^{2}+x+1$: $\Delta=1-4=-3$, $\delta=i\sqrt3$; $x=\frac{-1\pm i\sqrt3}{2}$.
  \item $x^{2}+x+2$: $\Delta=1-8=-7$, $\delta=i\sqrt7$; $x=\frac{-1\pm i\sqrt7}{2}$.
  \item $x^{2}+ix+1$: $\Delta=i^{2}-4=-5$, $\delta=i\sqrt5$; $x=\frac{-i\pm i\sqrt5}{2}=\frac{-1\pm\sqrt5}{2}\,i$.
\end{enumerate}
\resultado $\mathbf{\frac{-1\pm i\sqrt3}{2}}$; $\mathbf{\frac{-1\pm i\sqrt7}{2}}$; $\mathbf{\frac{-1+\sqrt5}{2}i\approx0{,}618\,i}$ \textbf{y} $\mathbf{\frac{-1-\sqrt5}{2}i\approx-1{,}618\,i}$.
\comprobacion en el tercero, suma $=-i=-b/a$ y producto $=i^{2}\cdot\frac{1-5}{4}=1=c/a$. En los dos primeros las raíces son conjugadas, como predice la \cref{prop:conjugadas}; en el tercero no, porque un coeficiente no es real.
\end{solucion}

\begin{solucion}{eje:5-suma-tripos}
\planteamiento agrupar las potencias de $i$ en bloques de cuatro, que suman cero.
\pasos
\begin{enumerate}[label=\arabic*.]
  \item $\sum_{a=0}^{200}i^{a}$: los términos $a=0,\dots,199$ forman $50$ bloques nulos; queda $i^{200}=1$.
  \item $\sum_{b=1}^{50}i^{b}$: los términos $b=1,\dots,48$ forman $12$ bloques nulos; quedan $i^{49}+i^{50}=i-1$.
  \item $x+iy=1+(i-1)=i$, luego $x=0$ e $y=1$.
\end{enumerate}
\resultado $\mathbf{xy=0}$.
\comprobacion suma geométrica: $\sum_{a=0}^{200}i^{a}=\frac{i^{201}-1}{i-1}=\frac{i-1}{i-1}=1$.
\end{solucion}

\begin{solucion}{eje:5-raices-cuadradas}
\planteamiento \cref{prop:raiz-cuadrada}.
\pasos
\begin{enumerate}[label=\arabic*.]
  \item $w=-5+12i$: $\abs{w}=13$; $x=\sqrt{\frac{13-5}{2}}=2$, $y=+\sqrt{\frac{13+5}{2}}=3$. Raíces $\pm(2+3i)$.
  \item $w=-2i$: $a=0$, $b=-2$, $\abs{w}=2$; $x=\sqrt1=1$, $y=-\sqrt1=-1$. Raíces $\pm(1-i)$.
\end{enumerate}
\resultado $\mathbf{\pm(2+3i)}$ \textbf{y} $\mathbf{\pm(1-i)}$.
\comprobacion $(2+3i)^{2}=4+12i-9=-5+12i$; $(1-i)^{2}=1-2i-1=-2i$.
\end{solucion}

\begin{solucion}{eje:5-cuadratica-compleja}
\datos $a=1$, $b=-(3+i)$, $c=2+2i$.
\planteamiento \cref{prop:cuadratica}, con la raíz cuadrada del discriminante del \cref{eje:5-raices-cuadradas}.
\pasos
\begin{enumerate}[label=\arabic*.]
  \item $b^{2}-4ac=(3+i)^{2}-4(2+2i)=8+6i-8-8i=-2i$.
  \item $\delta=1-i$ (\cref{eje:5-raices-cuadradas}).
  \item $z=\frac{3+i\pm(1-i)}{2}$: $z_1=\frac{4}{2}=2$, $z_2=\frac{2+2i}{2}=1+i$.
\end{enumerate}
\resultado $\mathbf{z\in\{2,\ 1+i\}}$.
\comprobacion $z_1+z_2=3+i=-b/a$ y $z_1z_2=2+2i=c/a$.
\end{solucion}

\begin{solucion}{eje:5-lugares}
\planteamiento escribir $z=x+iy$ y elevar al cuadrado las igualdades de módulos.
\pasos
\begin{enumerate}[label=\arabic*.]
  \item (a) $(x-1)^{2}+y^{2}=x^{2}+(y+1)^{2}$, es decir, $-2x+1=2y+1$: la recta $y=-x$, mediatriz del segmento que une $1$ y $-i$.
  \item (b) $x^{2}+(y+3)^{2}=9(x^{2}+y^{2})$, es decir, $8x^{2}+8y^{2}-6y-9=0$.
  \item Dividiendo por~8 y completando el cuadrado: $x^{2}+\bigl(y-\frac38\bigr)^{2}=\frac98+\frac{9}{64}=\frac{81}{64}$.
\end{enumerate}
\resultado \textbf{(a) la recta $y=-x$; (b) la circunferencia de centro $\frac38i$ y radio $\frac98$.}
\comprobacion los puntos $\frac32i$ y $-\frac34i$ de la circunferencia cumplen $\abs{\frac32i+3i}=\frac92=3\cdot\frac32$ y $\abs{-\frac34i+3i}=\frac94=3\cdot\frac34$.
\end{solucion}

\begin{solucion}{eje:5-escalar}
\datos $z=1+2i$, $w=3+i$.
\planteamiento \cref{prop:escalar-complejo} y \cref{teo:escalar}.
\pasos
\begin{enumerate}[label=\arabic*.]
  \item $z\bar w=(1+2i)(3-i)=3-i+6i-2i^{2}=5+5i$.
  \item $\Real(z\bar w)=5$ y $(1,2)\cdot(3,1)=3+2=5$.
  \item $\cos\varphi=\frac{5}{\sqrt5\sqrt{10}}=\frac{1}{\sqrt2}$, luego $\varphi=\frac{\pi}{4}$.
\end{enumerate}
\resultado \textbf{$z\bar w=5+5i$; producto escalar $5$; ángulo $\pi/4$.}
\comprobacion $\Arg z-\Arg w=\arctan2-\arctan\frac13=1{,}1071-0{,}3218=0{,}7854=\pi/4$.
\end{solucion}

\begin{solucion}{eje:5-modulo-producto}
\datos $z=a+bi$, $w=c+di$.
\planteamiento desarrollar en coordenadas.
\pasos
\begin{enumerate}[label=\arabic*.]
  \item $z\bar z=(a+bi)(a-bi)=a^{2}-abi+abi-b^{2}i^{2}=a^{2}+b^{2}$.
  \item $wz=(ac-bd)+(ad+bc)i$, luego $\abs{wz}^{2}=(ac-bd)^{2}+(ad+bc)^{2}$.
  \item Desarrollando: $a^{2}c^{2}-2abcd+b^{2}d^{2}+a^{2}d^{2}+2abcd+b^{2}c^{2}=(a^{2}+b^{2})(c^{2}+d^{2})=\abs{z}^{2}\abs{w}^{2}$.
  \item Se toman raíces no negativas.
\end{enumerate}
\resultado \textbf{$\abs{wz}=\abs{w}\abs{z}$, equivalente a la identidad real $(a^{2}+b^{2})(c^{2}+d^{2})=(ac-bd)^{2}+(ad+bc)^{2}$} (identidad de Brahmagupta--Fibonacci).
\comprobacion $(1+4)(9+1)=50$ y, con $z=1+2i$, $w=3+i$: $(3-2)^{2}+(1+6)^{2}=1+49=50$.
\end{solucion}

\begin{solucion}{eje:5-igualdad-triangular}
\datos $z,w\ne0$.
\planteamiento examinar cuándo es igualdad la única desigualdad de la demostración de la \cref{prop:modulo}(iv).
\pasos
\begin{enumerate}[label=\arabic*.]
  \item $\abs{z+w}^{2}=\abs{z}^{2}+\abs{w}^{2}+2\Real(z\bar w)$ y $(\abs{z}+\abs{w})^{2}=\abs{z}^{2}+\abs{w}^{2}+2\abs{z\bar w}$. Hay igualdad si y solo si $\Real(u)=\abs{u}$ con $u=z\bar w$.
  \item Para $u=p+qi$, $p\le\sqrt{p^{2}+q^{2}}$ con igualdad si y solo si $q=0$ y $p\ge0$: $u$ es real no negativo, y de hecho positivo, porque $u\ne0$.
  \item $z\bar w=t>0$ implica $\bar w=t/z$, luego $w=t/\bar z=tz/\abs{z}^{2}=\lambda z$ con $\lambda=t/\abs{z}^{2}>0$.
  \item Recíprocamente, si $w=\lambda z$ con $\lambda>0$: $\abs{z+\lambda z}=(1+\lambda)\abs{z}=\abs{z}+\abs{\lambda z}$.
\end{enumerate}
\resultado \textbf{$\abs{z+w}=\abs{z}+\abs{w}$ si y solo si $w=\lambda z$ con $\lambda>0$}: los dos vectores tienen la misma dirección y sentido.
\comprobacion $z=1+i$, $w=2+2i$: $\abs{3+3i}=3\sqrt2=\sqrt2+2\sqrt2$.
\end{solucion}

\begin{solucion}{eje:5-orden}
\datos un orden total $<$ en $\C$ compatible con la suma y el producto.
\planteamiento probar que todo cuadrado no nulo es positivo y aplicarlo a $1$ y a $i$.
\pasos
\begin{enumerate}[label=\arabic*.]
  \item Si $z\ne0$, por ser el orden total, $0<z$ o $z<0$. Si $0<z$, $0<z\cdot z$. Si $z<0$, sumando $-z$: $0<-z$, y $0<(-z)(-z)=z^{2}$. En ambos casos $0<z^{2}$.
  \item Con $z=1$: $0<1$. Con $z=i$: $0<i^{2}=-1$.
  \item Sumando $1$ a $0<-1$: $1<0$. Junto con $0<1$, la transitividad da $0<0$, que contradice que el orden es estricto.
\end{enumerate}
\resultado \textbf{No existe ningún orden en $\C$ compatible con sus operaciones.}
\comprobacion el argumento solo usa $i^{2}=-1$: falla en $\R$, donde $-1$ no es un cuadrado.
\end{solucion}

\begin{solucion}{eje:5-blaschke}
\datos $\abs{z}=1$, es decir, $z\bar z=1$; $\bar wz\ne1$.
\planteamiento escribir el denominador como un producto en el que aparezca $\overline{z-w}$.
\pasos
\begin{enumerate}[label=\arabic*.]
  \item $1-\bar wz=z\bar z-\bar wz=z(\bar z-\bar w)=z\,\overline{(z-w)}$.
  \item Por la \cref{prop:modulo}: $\abs{1-\bar wz}=\abs{z}\,\abs{\overline{z-w}}=1\cdot\abs{z-w}$.
  \item Como $\bar wz\ne1$, el denominador no se anula, y el cociente tiene módulo $\abs{z-w}/\abs{z-w}=1$ (si $z=w$ sería $\bar wz=\abs{z}^{2}=1$, excluido, luego $z-w\ne0$).
\end{enumerate}
\resultado $\mathbf{\left|\frac{z-w}{1-\bar wz}\right|=1}$.
\comprobacion $z=i$, $w=\frac12$: $\frac{i-\frac12}{1-\frac{i}{2}}$ tiene módulo $\frac{\sqrt{5}/2}{\sqrt{5}/2}=1$.
\end{solucion}

% ---------------------------------------------------------------------
\section{Tema 6}

\begin{solucion}{eje:6-exponenciales}
\planteamiento módulo y argumento de cada número (\cref{prop:arg}).
\pasos
\begin{enumerate}[label=\arabic*.]
  \item $1=e^{i0}$; \quad $i=e^{i\pi/2}$; \quad $-1=e^{i\pi}$.
  \item $1+i$: módulo $\sqrt2$, argumento $\frac{\pi}{4}$: $\sqrt2\,e^{i\pi/4}$.
  \item $\frac{1+\sqrt3i}{2}$: módulo $\sqrt{\frac14+\frac34}=1$, argumento $\arctan\sqrt3=\frac{\pi}{3}$: $e^{i\pi/3}$.
\end{enumerate}
\resultado $\mathbf{e^{0},\ e^{i\pi/2},\ e^{i\pi},\ \sqrt2e^{i\pi/4},\ e^{i\pi/3}}$.
\comprobacion $\sqrt2e^{i\pi/4}=\sqrt2(\frac{\sqrt2}{2}+\frac{\sqrt2}{2}i)=1+i$.
\end{solucion}

\begin{solucion}{eje:6-vf}
\pasos
\begin{enumerate}[label=(\alph*)]
  \item \textbf{Verdadera.} $\cos\pi+i\sin\pi=-1$.
  \item \textbf{Verdadera.} \Cref{prop:exp}(ii).
  \item \textbf{Falsa.} $z=w=-1$: $\Arg(zw)=\Arg1=0$, pero $\Arg z+\Arg w=2\pi$.
  \item \textbf{Falsa.} $e^{i0}=e^{2\pi i}=1$; solo se deduce $\theta-\varphi\in2\pi\Z$.
  \item \textbf{Falsa.} Un giro que fija el punto $1$ es la identidad (de $e^{i\alpha}\cdot1=1$ sale $e^{i\alpha}=1$), pero la conjugación fija $1$ y lleva $i$ a $-i$. Es la reflexión en el eje real.
\end{enumerate}
\resultado \textbf{(a) y (b) verdaderas; (c), (d) y (e) falsas.}
\end{solucion}

\begin{solucion}{eje:6-cuadrante}
\planteamiento los argumentos se suman al multiplicar y se restan al dividir (\cref{teo:producto-polar}).
\pasos
\begin{enumerate}[label=\arabic*.]
  \item $\arg(1+i)=\frac{\pi}{4}$, luego $\arg z=5\cdot\frac{\pi}{4}=\frac{5\pi}{4}$: cuadrante III.
  \item $\arg(1-i)=-\frac{\pi}{4}$, luego $\arg w=\frac{\pi}{4}-2\cdot(-\frac{\pi}{4})=\frac{3\pi}{4}$: cuadrante II.
\end{enumerate}
\resultado \textbf{$z$ en el cuadrante III; $w$ en el cuadrante II.}
\comprobacion $(1+i)^{5}=(1+i)^{4}(1+i)=-4-4i$; $(1-i)^{2}=-2i$ y $w=\frac{1+i}{-2i}=-\frac12+\frac12i$.
\end{solucion}

\begin{solucion}{eje:6-polar}
\planteamiento módulo y \cref{prop:arg}.
\pasos
\begin{enumerate}[label=\arabic*.]
  \item $-1+i$: $r=\sqrt2$; $a<0\le b$: $\arctan(-1)+\pi=\frac{3\pi}{4}$.
  \item $-\sqrt3-i$: $r=2$; $a<0$, $b<0$: $\arctan\frac{1}{\sqrt3}-\pi=\frac{\pi}{6}-\pi=-\frac{5\pi}{6}$.
  \item $-2i$: $r=2$; $a=0$, $b<0$: $-\frac{\pi}{2}$.
  \item $3-4i$: $r=5$; $a>0$: $\arctan(-\frac43)\approx-0{,}9273$.
\end{enumerate}
\resultado $\mathbf{\sqrt2e^{3\pi i/4},\ 2e^{-5\pi i/6},\ 2e^{-\pi i/2},\ 5e^{-0{,}9273\,i}}$.
\comprobacion $2e^{-5\pi i/6}=2(-\frac{\sqrt3}{2}-\frac12i)=-\sqrt3-i$.
\end{solucion}

\begin{solucion}{eje:6-producto-polar}
\planteamiento \cref{teo:producto-polar}: se multiplican o dividen los módulos y se suman o restan los argumentos.
\pasos
\begin{enumerate}[label=\arabic*.]
  \item $z_1z_2=60\,e^{i(\pi/6+\pi/3)}=60e^{i\pi/2}=60i$.
  \item $z_1/z_2=\frac{10}{6}e^{i(\pi/6-\pi/3)}=\frac53e^{-i\pi/6}=\frac53\bigl(\frac{\sqrt3}{2}-\frac12i\bigr)=\frac{5\sqrt3}{6}-\frac56i$.
  \item $12e^{i55^{\circ}}\cdot3e^{i18^{\circ}}=36e^{i73^{\circ}}$; $\;12e^{i55^{\circ}}/3e^{i18^{\circ}}=4e^{i37^{\circ}}$.
\end{enumerate}
\resultado $\mathbf{60i}$; $\mathbf{\frac53e^{-i\pi/6}\approx1{,}443-0{,}8333\,i}$; $\mathbf{36e^{i73^{\circ}}}$; $\mathbf{4e^{i37^{\circ}}}$.
\comprobacion en binómica, $z_1=5\sqrt3+5i$, $z_2=3+3\sqrt3i$, y $z_1z_2=15\sqrt3+45i+15i-15\sqrt3=60i$.
\end{solucion}

\begin{solucion}{eje:6-potencias}
\planteamiento \cref{teo:demoivre}.
\pasos
\begin{enumerate}[label=\arabic*.]
  \item $\frac{1+\sqrt3i}{2}=e^{i\pi/3}$, luego su cubo es $e^{i\pi}=-1$.
  \item $1+i=\sqrt2e^{i\pi/4}$, luego $(1+i)^{4}=4e^{i\pi}=-4$.
  \item En binómica, $\bigl(\frac{1+\sqrt3i}{2}\bigr)^{2}=\frac{1+2\sqrt3i-3}{4}=\frac{-1+\sqrt3i}{2}$. Multiplicando de nuevo por $\frac{1+\sqrt3i}{2}$ se obtiene $\frac{(-1+\sqrt3i)(1+\sqrt3i)}{4}=\frac{-1-3}{4}=-1$.
\end{enumerate}
\resultado $\mathbf{-1}$ \textbf{y} $\mathbf{-4}$.
\comprobacion $(1+i)^{4}=-4$ también en el \cref{eje:5-operaciones}.
\end{solucion}

\begin{solucion}{eje:6-tripos}
\planteamiento $\frac{1+i}{1-i}=i$; separar el cociente en una potencia de~$i$ y un factor sobrante.
\pasos
\begin{enumerate}[label=\arabic*.]
  \item $\frac{1+i}{1-i}=\frac{(1+i)^{2}}{2}=\frac{2i}{2}=i$.
  \item $\frac{(1+i)^{2019}}{(1-i)^{2017}}=\Bigl(\frac{1+i}{1-i}\Bigr)^{2017}(1+i)^{2}=i^{2017}\cdot2i$.
  \item $2017=4\cdot504+1$: $i^{2017}=i$, luego el resultado es $2i^{2}=-2$.
\end{enumerate}
\resultado $\mathbf{-2}$.
\comprobacion en forma polar: $\frac{2^{2019/2}e^{2019\pi i/4}}{2^{2017/2}e^{-2017\pi i/4}}=2e^{4036\pi i/4}=2e^{1009\pi i}=-2$.
\end{solucion}

\begin{solucion}{eje:6-raices}
\planteamiento \cref{teo:raices}.
\pasos
\begin{enumerate}[label=\arabic*.]
  \item (a) $e^{ik\pi/4}$, $k=0,\dots,7$: $1$, $\frac{\sqrt2}{2}(1+i)$, $i$, $\frac{\sqrt2}{2}(-1+i)$, $-1$, $\frac{\sqrt2}{2}(-1-i)$, $-i$, $\frac{\sqrt2}{2}(1-i)$.
  \item (b) $-8=8e^{i\pi}$; raíces $2e^{i(\pi+2\pi k)/3}$: $2e^{i\pi/3}=1+\sqrt3i$, $2e^{i\pi}=-2$, $2e^{5\pi i/3}=1-\sqrt3i$. Son los vértices de un triángulo equilátero inscrito en la circunferencia de radio~2.
  \item (c) $-\frac{i}{64}=\frac1{64}e^{-i\pi/2}$; $(\frac1{64})^{1/6}=\frac12$; argumentos $\frac{-\pi/2+2\pi k}{6}=-\frac{\pi}{12}+\frac{k\pi}{3}$.
\end{enumerate}
\resultado \textbf{(a) las ocho anteriores; (b) $1\pm\sqrt3i$ y $-2$; (c) $\frac12e^{i(-\pi/12+k\pi/3)}$, $k=0,\dots,5$}, es decir, argumentos $-15^{\circ}$, $45^{\circ}$, $105^{\circ}$, $165^{\circ}$, $225^{\circ}$ y $285^{\circ}$.
\comprobacion $(1+\sqrt3i)^{3}=(2e^{i\pi/3})^{3}=8e^{i\pi}=-8$; $(\frac12e^{-i\pi/12})^{6}=\frac1{64}e^{-i\pi/2}=-\frac{i}{64}$.
\end{solucion}

\begin{solucion}{eje:6-polinomios}
\planteamiento raíces $n$-ésimas y factorización directa.
\pasos
\begin{enumerate}[label=\arabic*.]
  \item $x^{3}=-1=e^{i\pi}$: $x=e^{i\pi/3},\,e^{i\pi},\,e^{5i\pi/3}$, es decir, $\frac{1\pm\sqrt3i}{2}$ y $-1$.
  \item $x^{4}-2x^{2}+1=(x^{2}-1)^{2}$: raíces $1$ y $-1$, ambas dobles.
  \item $x^{4}=-4=4e^{i\pi}$: $x=\sqrt2e^{i(\pi+2\pi k)/4}$, es decir, $1+i$, $-1+i$, $-1-i$, $1-i$.
  \item Agrupando conjugadas: $(x-1-i)(x-1+i)=x^{2}-2x+2$ y $(x+1-i)(x+1+i)=x^{2}+2x+2$.
\end{enumerate}
\resultado \textbf{$-1$, $\frac{1\pm\sqrt3i}{2}$; $\pm1$ dobles; $\pm1\pm i$; $x^{4}+4=(x^{2}-2x+2)(x^{2}+2x+2)$.}
\comprobacion $(x^{2}+2)^{2}-(2x)^{2}=x^{4}+4x^{2}+4-4x^{2}=x^{4}+4$.
\end{solucion}

\begin{solucion}{eje:6-giros}
\planteamiento interpretación geométrica del producto y \cref{prop:giro-complejo}.
\pasos
\begin{enumerate}[label=\arabic*.]
  \item (a) $i=e^{i\pi/2}$: giro de $90^{\circ}$. $-1=e^{i\pi}$: media vuelta (simetría central). $2i$: giro de $90^{\circ}$ seguido de homotecia de razón~2. $\frac{1+i}{\sqrt2}=e^{i\pi/4}$: giro de $45^{\circ}$.
  \item (b) $(2+i)\,i=-1+2i$, es decir, el punto $(-1,2)$; con la matriz, $R(\frac{\pi}{2})(2,1)\T=\begin{pmatrix}0&-1\\1&0\end{pmatrix}\begin{pmatrix}2\\1\end{pmatrix}=\begin{pmatrix}-1\\2\end{pmatrix}$.
  \item $(2+i)e^{i\pi/3}=(2+i)\bigl(\frac12+\frac{\sqrt3}{2}i\bigr)=\bigl(1-\frac{\sqrt3}{2}\bigr)+\bigl(\sqrt3+\frac12\bigr)i\approx0{,}1340+2{,}232\,i$.
  \item (c) $c+i(z-c)=(1+i)+i\cdot2=1+3i$.
\end{enumerate}
\resultado \textbf{(b) $(-1,2)$ y $\approx(0{,}1340;\,2{,}232)$; (c) $1+3i$.}
\comprobacion los giros conservan la distancia al centro: $\abs{0{,}1340+2{,}232i}^{2}=0{,}018+4{,}982=5=\abs{2+i}^{2}$, y $\abs{(1+3i)-c}=\abs{2i}=2=\abs{(3+i)-c}$.
\end{solucion}

\begin{solucion}{eje:6-rope}
\datos $z_1=1+2i$, $z_2=3+i$, $\theta=\pi/6$, $(m,n)=(3,1)$ y $(5,3)$.
\planteamiento calcular directamente y con el \cref{cor:rope}.
\pasos
\begin{enumerate}[label=\arabic*.]
  \item Directo: $z_1e^{i\pi/2}=i(1+2i)=-2+i$; $z_2e^{i\pi/6}=(3+i)\bigl(\frac{\sqrt3}{2}+\frac12i\bigr)=\frac{3\sqrt3-1}{2}+\frac{3+\sqrt3}{2}i\approx2{,}098+2{,}366i$.
  \item Producto escalar de $(-2,1)$ y $(2{,}098;\,2{,}366)$: $-4{,}196+2{,}366=-1{,}830$.
  \item Con el corolario: $z_1\bar z_2=5+5i$ (\cref{eje:5-escalar}) y $e^{i(m-n)\theta}=e^{i\pi/3}=\frac12+\frac{\sqrt3}{2}i$; $\Real\bigl((5+5i)(\frac12+\frac{\sqrt3}{2}i)\bigr)=\frac52-\frac{5\sqrt3}{2}$.
  \item Con $(m,n)=(5,3)$ la diferencia es también~2, y el corolario da el mismo valor.
\end{enumerate}
\resultado $\mathbf{\frac{5-5\sqrt3}{2}\approx-1{,}830}$ \textbf{en ambos casos.}
\comprobacion coincide con el ejemplo de RoPE del módulo~14 (tema~4), con $q=(1,2)$, $k=(3,1)$ y distancia $n-m=-2$.
\end{solucion}

\begin{solucion}{eje:6-suma-euler}
\datos $e^{i\theta}=\cos\theta+i\sin\theta$; $e^{z+w}=e^{z}e^{w}$.
\planteamiento desarrollar $e^{i(\alpha+\beta)}=e^{i\alpha}e^{i\beta}$ e igualar partes reales e imaginarias.
\pasos
\begin{enumerate}[label=\arabic*.]
  \item Miembro izquierdo: $\cos(\alpha+\beta)+i\sin(\alpha+\beta)$.
  \item Miembro derecho: $(\cos\alpha+i\sin\alpha)(\cos\beta+i\sin\beta)=(\cos\alpha\cos\beta-\sin\alpha\sin\beta)+i(\sin\alpha\cos\beta+\cos\alpha\sin\beta)$.
  \item Dos complejos son iguales si y solo si coinciden sus partes reales e imaginarias, que aquí son números reales.
\end{enumerate}
\resultado $\mathbf{\cos(\alpha+\beta)=\cos\alpha\cos\beta-\sin\alpha\sin\beta}$ \textbf{y} $\mathbf{\sin(\alpha+\beta)=\sin\alpha\cos\beta+\cos\alpha\sin\beta}$.
\comprobacion con $\beta=\alpha$ se obtienen las fórmulas del ángulo doble de la \cref{prop:doble}.
\end{solucion}

\begin{solucion}{eje:6-cos5}
\datos $c=\cos\theta$, $s=\sin\theta$; $(c+is)^{5}=\cos5\theta+i\sin5\theta$.
\planteamiento binomio de Newton y parte real; después, $\theta=\pi/10$.
\pasos
\begin{enumerate}[label=\arabic*.]
  \item $(c+is)^{5}=\sum_{k=0}^{5}\binom5kc^{5-k}(is)^{k}$; los términos reales son los de $k=0,2,4$: $c^{5}-10c^{3}s^{2}+5cs^{4}$.
  \item Con $s^{2}=1-c^{2}$: $c^{5}-10c^{3}+10c^{5}+5c(1-2c^{2}+c^{4})=16c^{5}-20c^{3}+5c$.
  \item (b) $\theta=\frac{\pi}{10}$: $\cos\frac{\pi}{2}=0$, y como $c>0$, $16c^{4}-20c^{2}+5=0$.
  \item $c^{2}=\frac{20\pm\sqrt{400-320}}{32}=\frac{5\pm\sqrt5}{8}$.
  \item Como $0<\frac{\pi}{10}<\frac{\pi}{4}$ y el coseno es decreciente en $[0,\frac{\pi}{2}]$, $c^{2}>\cos^{2}\frac{\pi}{4}=\frac12$; solo $\frac{5+\sqrt5}{8}\approx0{,}9045$ lo cumple ($\frac{5-\sqrt5}{8}\approx0{,}3455$).
\end{enumerate}
\resultado $\mathbf{\cos5\theta=16\cos^{5}\theta-20\cos^{3}\theta+5\cos\theta}$ \textbf{y} $\mathbf{\cos^{2}\frac{\pi}{10}=\frac{5+\sqrt5}{8}}$.
\comprobacion $\cos^{2}18^{\circ}=0{,}951\,06^{2}=0{,}9045$. La otra raíz es $\cos^{2}\frac{3\pi}{10}$, porque $\cos\frac{15\pi}{10}=0$ también.
\end{solucion}

\begin{solucion}{eje:6-potencias-raices}
\datos $\omega=e^{2\pi i/n}$, $n\ge2$, $j\in\Z$.
\planteamiento suma geométrica de razón $\omega^{j}$; distinguir si la razón es~1.
\pasos
\begin{enumerate}[label=\arabic*.]
  \item $\omega^{j}=e^{2\pi ij/n}=1$ si y solo si $j/n\in\Z$ (\cref{prop:exp}(iii)), es decir, si $n$ divide a $j$. En ese caso la suma es $n\cdot1=n$.
  \item Si $n$ no divide a $j$: $\sum_{k=0}^{n-1}(\omega^{j})^{k}=\frac{(\omega^{j})^{n}-1}{\omega^{j}-1}=\frac{e^{2\pi ij}-1}{\omega^{j}-1}=0$.
  \item Con $j=1$ (que $n\ge2$ no divide): $\sum_{k}e^{2\pi ik/n}=0$; las partes real e imaginaria son las sumas de cosenos y senos pedidas.
\end{enumerate}
\resultado \textbf{La suma vale $n$ si $n\mid j$ y $0$ en otro caso; $\sum_k\cos\frac{2\pi k}{n}=\sum_k\sin\frac{2\pi k}{n}=0$.}
\comprobacion $n=3$: $\cos0+\cos\frac{2\pi}{3}+\cos\frac{4\pi}{3}=1-\frac12-\frac12=0$.
\end{solucion}

\begin{solucion}{eje:6-suma-cosenos}
\datos $x\notin2\pi\Z$, de modo que $e^{ix}\ne1$ y $\sin\frac x2\ne0$.
\planteamiento suma geométrica de razón $e^{ix}$ y factorización $e^{i\varphi}-1=e^{i\varphi/2}\,2i\sin\frac{\varphi}{2}$.
\pasos
\begin{enumerate}[label=\arabic*.]
  \item La indicación: $e^{i\varphi/2}(e^{i\varphi/2}-e^{-i\varphi/2})=e^{i\varphi/2}\cdot2i\sin\frac{\varphi}{2}$, por la \cref{prop:exp}(iv).
  \item $\sum_{k=0}^{n}e^{ikx}=\frac{e^{i(n+1)x}-1}{e^{ix}-1}=\frac{e^{i(n+1)x/2}\,2i\sin\frac{(n+1)x}{2}}{e^{ix/2}\,2i\sin\frac{x}{2}}=e^{inx/2}\,\frac{\sin\frac{(n+1)x}{2}}{\sin\frac x2}$.
  \item La parte real del miembro izquierdo es $\sum\cos kx$, y la del derecho, $\cos\frac{nx}{2}\cdot\frac{\sin\frac{(n+1)x}{2}}{\sin\frac x2}$.
\end{enumerate}
\resultado $\displaystyle\sum_{k=0}^{n}\cos kx=\frac{\cos\frac{nx}{2}\sin\frac{(n+1)x}{2}}{\sin\frac x2}$.
\comprobacion la parte imaginaria da $\sum_{k=1}^{n}\sin kx=\frac{\sin\frac{nx}{2}\sin\frac{(n+1)x}{2}}{\sin\frac x2}$, el resultado del \cref{eje:2-suma-senos}. Con $n=7$ y $x=0{,}9$, ambos miembros valen $1{,}0173$.
\end{solucion}

\begin{solucion}{eje:6-autovalores}
\datos $R(\alpha)$ con $\sin\alpha\ne0$.
\planteamiento polinomio característico y comprobación directa de los autovectores.
\pasos
\begin{enumerate}[label=\arabic*.]
  \item $\det(R(\alpha)-\lambda I)=(\cos\alpha-\lambda)^{2}+\sin^{2}\alpha=\lambda^{2}-2\lambda\cos\alpha+1$.
  \item Discriminante: $4\cos^{2}\alpha-4=-4\sin^{2}\alpha<0$: no hay raíces reales.
  \item Raíces complejas: $\lambda=\cos\alpha\pm i\sin\alpha=e^{\pm i\alpha}$.
  \item $R(\alpha)\begin{pmatrix}1\\-i\end{pmatrix}=\begin{pmatrix}\cos\alpha+i\sin\alpha\\\sin\alpha-i\cos\alpha\end{pmatrix}=e^{i\alpha}\begin{pmatrix}1\\-i\end{pmatrix}$, porque $e^{i\alpha}(-i)=\sin\alpha-i\cos\alpha$. Del mismo modo, $R(\alpha)(1,i)\T=e^{-i\alpha}(1,i)\T$.
\end{enumerate}
\resultado \textbf{Autovalores $e^{\pm i\alpha}$, no reales, con autovectores $(1,\mp i)\T$.} Por la \cref{prop:giro-complejo}(i), $R(\alpha)$ actúa sobre el plano como la multiplicación por $e^{i\alpha}$, y ese número aparece como autovalor.
\comprobacion producto de autovalores $e^{i\alpha}e^{-i\alpha}=1=\det R(\alpha)$; suma $2\cos\alpha=\operatorname{tr}R(\alpha)$.
\end{solucion}
